% =====================================================================
%  PoP -> IEEE TDSC : 英文论文骨架
%  生成日期 2026-10-12 · 配套 docs/tdsc-plan.md
%
%  用法：这是 SKELETON，不是成稿。所有 \need{} 与 \wpmark{} 处都要填。
%  决策锁定：见 docs/tdsc-plan.md §0（D1–D10）。
%
%  编译校验（本机已验通）：
%      pdflatex "\def\SKELQUIET{}% =====================================================================
%  PoP -> IEEE TDSC : 英文论文骨架
%  生成日期 2026-10-12 · 配套 docs/tdsc-plan.md
%
%  用法：这是 SKELETON，不是成稿。所有 \need{} 与 \wpmark{} 处都要填。
%  决策锁定：见 docs/tdsc-plan.md §0（D1–D10）。
%
%  编译校验（本机已验通）：
%      pdflatex "\def\SKELQUIET{}% =====================================================================
%  PoP -> IEEE TDSC : 英文论文骨架
%  生成日期 2026-10-12 · 配套 docs/tdsc-plan.md
%
%  用法：这是 SKELETON，不是成稿。所有 \need{} 与 \wpmark{} 处都要填。
%  决策锁定：见 docs/tdsc-plan.md §0（D1–D10）。
%
%  编译校验（本机已验通）：
%      pdflatex "\def\SKELQUIET{}% =====================================================================
%  PoP -> IEEE TDSC : 英文论文骨架
%  生成日期 2026-10-12 · 配套 docs/tdsc-plan.md
%
%  用法：这是 SKELETON，不是成稿。所有 \need{} 与 \wpmark{} 处都要填。
%  决策锁定：见 docs/tdsc-plan.md §0（D1–D10）。
%
%  编译校验（本机已验通）：
%      pdflatex "\def\SKELQUIET{}\input{tdsc-skeleton.tex}"   % 两次，第二遍收交叉引用
%  \SKELQUIET 让所有中文标记宏吞掉参数；不带它则需要能排版中文的引擎
%  （xelatex + CJK 字体），本骨架未配置。
% =====================================================================
\documentclass[journal]{IEEEtran}

\usepackage{amsmath,amssymb,amsthm}
\usepackage{booktabs}
\usepackage{array}
\usepackage{graphicx}
\usepackage{xcolor}
\usepackage{hyperref}

% 窄栏里用 p{} 会因两端对齐产生大片字距（编译时刷一堆 Underfull \hbox）。
% 表格正文列一律用 L{}：左对齐、允许自然参差。
\newcolumntype{L}[1]{>{\raggedright\arraybackslash}p{#1}}
\newcommand{\tabfont}{\footnotesize}

\newtheorem{theorem}{Theorem}
\newtheorem{lemma}[theorem]{Lemma}
\newtheorem{proposition}[theorem]{Proposition}
\newtheorem{corollary}[theorem]{Corollary}
\newtheorem{definition}[theorem]{Definition}
\newtheorem{assumption}[theorem]{Assumption}
\newtheorem{remark}[theorem]{Remark}

% --- 工作标记：提交前全部清掉 ---
% 注意：宏名不能叫 \wp —— \wp 已是 LaTeX/amssymb 的 Weierstrass ℘，
% \newcommand{\wp} 会直接报 "Command \wp already defined" 并使后文级联报错。
\ifdefined\SKELQUIET
  \newcommand{\wpmark}[1]{}
  \newcommand{\need}[1]{}
  \newcommand{\fixme}[1]{}
\else
  \newcommand{\wpmark}[1]{\textcolor{red}{\textsuperscript{[WP#1]}}}
  \newcommand{\need}[1]{\textcolor{blue}{[[ NEED: #1 ]]}}
  \newcommand{\fixme}[1]{\textcolor{orange}{[[ FIX: #1 ]]}}
\fi

\begin{document}

% ---------------------------------------------------------------------
% 标题：决策 D5 —— 保留 "Zero-Knowledge" 的前提是 groth16 成为交付档 (WP1)。
%       若 WP1 未打通 groth16，必须删掉标题与摘要里的 zero-knowledge。
% ---------------------------------------------------------------------
\title{Proof-of-Policy: Zero-Knowledge Compliance Certificates\\ for LLM Agent Outputs}
% 备选标题（若 groth16 未打通）："Proof-of-Policy: Succinct Compliance
% Certificates for LLM Agent Outputs"

\author{\need{authors}}
\maketitle

% ---------------------------------------------------------------------
\begin{abstract}
% 决策 E2：内容隐私上界必须出现在摘要，不能埋进 §8.1。
We present \emph{Proof-of-Policy} (PoP), a system that proves an LLM agent's
response $T$ satisfies a publicly declared policy $\pi$, without a trusted
execution environment...
\need{一句完整的贡献陈述（对应 D1 的"分析 x 系统各半"）}

\smallskip
\noindent\textbf{Scope of privacy.} PoP provides \emph{policy}-zero-knowledge:
the certificate reveals neither the policy evaluation trace nor, in
policy-private mode, the policy itself. It does \emph{not} guarantee that $T$
is unrecoverable: as we show in \S\ref{sec:binding}, any publicly recomputable
binding to $T$ admits a re-execution attack whenever the adversary knows the
prompt and the model. The same bit-budget governs the request side: a
certificate is replayable without limit unless the receiver contributes fresh
randomness (Theorem \ref{thm:t3}).
\end{abstract}

\begin{IEEEkeywords}
zero-knowledge proofs, LLM agents, policy compliance, verifiable computation,
accountability
\end{IEEEkeywords}

% =====================================================================
\section{Introduction}
\label{sec:intro}

\need{Contribution list. 必须包含下面这条"不主张"声明：}

\noindent\textbf{What we do \emph{not} claim.}
We do not claim that policy proofs bind nothing: our own system already binds
the response via a verifier-issued nonce (Definition 1). We do not claim to be
the first to observe that a proof can fail to bind the artifact it purports to
describe: the \emph{effort gap} for ZK inference proofs is formalised by
Hollow-LLM \cite{hollowllm}, and the \emph{source-integrity gap} for predicate
proofs is named by zk-attach \cite{zkattach}. Our contribution is to
(i) show that the surviving gap is the \emph{divergent-delivery} one
(Proposition \ref{thm:gap}), (ii) give a quantitative bound on what a public
value reveals, together with the \emph{structural hypothesis} under which that
bound holds---and a demonstration that the hypothesis is necessary, since our
own two public values sit on opposite sides of it (Theorem \ref{thm:t1} and
Remark \ref{rem:betanecessary}), and (iii) give a quantitative price for
closing the gap (Theorem \ref{thm:t3}).

\smallskip
\noindent\textbf{A separation we deliberately do \emph{not} claim.} It is
tempting to promote our designated-verifier interface into a theorem
(``public verification is impossible when $k$ is small''). We give the reason
we do not in Remark \ref{rem:designated}: in the low-$k$ regime hiding fails
for every scheme, designated or not, and the dichotomy has no content of its
own. We state this explicitly because the negated claim is the more attractive
one to write.

\smallskip
\noindent\textbf{One spine, two sides.} Theorems \ref{thm:t1} and
\ref{thm:t3} state the same fact about two different inputs: \emph{a proof
binds only as much as the committed, adversary-non-selectable quantity has bits
for}. The response side is governed by entropy,
$k = -\log \mathsf{Guess}(T \mid Z)$ (Definition \ref{def:k}); the request side
by freshness, $f$. Neither can be bought with more computation.

\fixme{全文禁止出现"比 Proof-of-Guardrail 更强"或任何强弱对比。
Proof-of-Guardrail 与本文正交：它证"计算发生过"，本文证"性质成立"。}

% =====================================================================
\section{Background and Related Work}
\label{sec:related}

\need{五条线，每条都要"它占什么位 / 我们差在哪"两列。见 docs/tdsc-plan.md \S2。}

\begin{table}[t]
\centering
\caption{Positioning against the five prior-art lineages.}
\label{tab:related}
\setlength{\tabcolsep}{3pt}
\begin{tabular}{@{}L{0.26\columnwidth}L{0.42\columnwidth}L{0.26\columnwidth}@{}}
\toprule
Lineage & Representative work & Occupied slot \\
\midrule
Policy $\to$ circuit      & zkGen \cite{zkgen}, Circuitree, Prezta & DSL-to-circuit \\
Agent/LLM compliance      & NiyamAI, Zaguk, PCO \cite{pco}, WITNESS & private-mode kernel \\
Model computation         & DeepProve, zkGPT, zkAgent \cite{zkagent} & inference integrity \\
Attacks / negative results& Hollow-LLM \cite{hollowllm}, zk-attach \cite{zkattach} & naming the gap \\
Non-ZK industry baseline  & @brivora/verify, VAOL & signed audit logs \\
\bottomrule
\end{tabular}
\end{table}

\fixme{引文纪律：除 Hollow-LLM / zkGen / NiyamAI / zkAgent / ePrint 2024/813
五条外，其余引文必须逐条打开原文核实后再进 bibliography。
见 docs/tdsc-fixlist.md \S2 的注记。}

% =====================================================================
\section{System Overview}
\label{sec:system}
% 决策 D1 的"系统半边"：这一节压缩到 ~1.5 页，是 T3(b) 的载体，不是主角。

\need{DSL $\to$ ConstraintSpec $\to$ circuit；双模式；证书 + 锚定；
工具回执链；801 测试 / 15 skip 的工件陈述。}

% =====================================================================
\section{Binding of Policy Proofs}
\label{sec:binding}
% 决策 D1 的"分析半边"：这是论文的核心。

\subsection{The gap}

\begin{definition}[Policy-proof system]
A policy-proof system is a tuple $(\mathsf{Prove}, \mathsf{Ver}, \pi, \mathcal{M})$.
On challenge $n$ from a verifier, $\mathsf{Prove}(\pi, T, n)$ outputs a
certificate $c$ whose public values include
\begin{multline*}
b \;=\; H(\texttt{pop-bind-v1} \,\|\, \mathsf{u32be}(|n|) \,\|\, n \,\|\, T),\\
C = H(T).
\end{multline*}
$\mathsf{Ver}(c, n, T')$ accepts iff the policy relation holds \emph{and}
$b$ recomputed from the received $T'$ matches.
\end{definition}

\fixme{\textbf{H3/H4}：早先写的 ``the relation contains no term referring to the
delivered bytes'' 与仓库事实冲突——公开值里\emph{已有} $b$，验证时由送达的
$T'$ 现场重算（\texttt{circuits/types/src/lib.rs:444}、
\texttt{docs/security-model.md:232}）。命题已按上面的定义改写。见
\texttt{docs/tdsc-fixlist.md} \S H。}

\begin{proposition}[Divergent delivery]
\label{thm:gap}
Definition 1 binds $T'$ \emph{as received by the verifying party}. If the
provider delivers a different string $T'' \ne T'$ to the end user, verification
still succeeds on $T'$. Neither hardware nor a proof can prevent this; only the
end user holding $n$ and recomputing $b$ can.
\end{proposition}
\begin{proof}
$b$ is a function of $(n,T)$ only; the delivery channel to the end user carries
no authenticated term. One line.
\end{proof}

\need{This is the \emph{residual} gap, not ``proofs bind nothing''. Prior work
must be cited here (Hollow-LLM \cite{hollowllm}, zk-attach \cite{zkattach}), and
we state plainly that the naive claim would be false against our own system.}

\subsection{A quantitative binding--hiding tradeoff}
\label{sec:tradeoff}

% ---------------------------------------------------------------------
% T1 : 关键是把"对手已知信息 Z"定对。
%      docs/tdsc-plan.md §1.4 已修正原 §4 的非因果推理。
% ---------------------------------------------------------------------
\begin{definition}[Guessing entropy of a response]
\label{def:k}
For a response distribution $\mathcal{D}$ and adversary side information $Z$,
let
\[
\mathsf{Guess}(T \mid Z) \;\triangleq\; \max_{\mathcal{A}} \Pr[\mathcal{A}(Z) = T],
\]
which equals $\mathbb{E}_z[\max_t \Pr[T = t \mid Z = z]]$, and
\[
k \;\triangleq\; -\log_2 \mathsf{Guess}(T \mid Z),
\]
the maximum ranging over all time-unbounded adversaries.
\end{definition}

\begin{remark}[The definition is operational, and it is the average]
\label{rem:kmeas}
$k$ is \emph{measured}, not assumed: sampling one model on one prompt many
times estimates $\mathsf{Guess}(T \mid Z)$ directly
(\S\ref{sec:eval}), so Theorem \ref{thm:t1} is instantiated with a
number we obtain rather than a number we posit. The definition must average
over $z$ rather than take the pointwise worst case
$k_z = -\log \max_t \Pr[T = t \mid Z = z]$: the proof takes an expectation over
$z$, and $\mathbb{E}_z[2^{-k_z}] = 2^{-k}$ holds only for the average form.
\end{remark}

\begin{assumption}[Adversary side information]
\label{as:Z}
The adversary knows the prompt, the model (possibly its weights), and the
decoding parameters. \emph{This is the assumption that makes the theorem
non-trivial}: under greedy decoding, $k \approx 0$.
\end{assumption}

\begin{theorem}[Binding, and the hypothesis it needs]
\label{thm:t1}
Let $\beta : \mathcal{M} \to \{0,1\}^\lambda$ be public, deterministic and
efficiently computable, and let the certificate publish $\beta(T)$. Let
$\mathcal{A}$ be any adversary, unbounded in time, that is given $Z$ and
evaluates $\beta$ at most $q$ times.
\begin{enumerate}
\item[(H)] \emph{(Randomised public value.)} Suppose
  $\beta(T) = H(\mathsf{dom} \,\|\, \mathsf{u32be}(|n|) \,\|\, n \,\|\, T)$ for a
  random oracle $H$ and a nonce $n$ drawn uniformly and \emph{independently of
  $T$}. Then
  \[
  \Pr\!\left[\mathcal{A}^\beta\!\left(Z, \beta(T)\right) = T\right]
  \;\le\; q \cdot 2^{-k} \;+\; q \cdot 2^{-\lambda}.
  \]
  Consequently hiding with advantage $\varepsilon$ requires
  $k \ge \log q + \log(1/\varepsilon)$.
\item[(L)] \emph{(Lossy public value.)} For an arbitrary deterministic public
  $\beta$ the correct bound is over the \emph{preimage class}:
  \begin{multline*}
  \Pr\!\left[\mathcal{A}^\beta\!\left(Z, \beta(T)\right) = T\right]
  \;\le\; \mathsf{Guess}\!\left(T \mid Z, \beta(T)\right) \\
  \triangleq\; \mathbb{E}_y\!\left[\max_t \Pr[T = t \mid Z, \beta(T) = y]\right].
  \end{multline*}
\end{enumerate}
\end{theorem}

\begin{proof}
\emph{(H).} Fix $z$ and let the adversary's query set be
$Q = \{q_1, \dots, q_q\}$. Decompose by first hit,
\[
\Pr[\exists i : q_i = T] \;=\; \sum_i \Pr\!\left[q_i = T \wedge \neg E_{<i}\right],
\]
where $E_{<i} \triangleq \{\exists j < i : q_j = T\}$.
Conditioned on $\neg E_{<i}$, every oracle answer $a_1,\dots,a_{i-1}$ already
revealed is uniform and independent of $T$ (no queried point equals $T$, and $H$
is a fresh random function), and $\beta(T)$ itself is uniform and independent of
$T$. Hence $q_i$---a deterministic function of $(z, \beta(T), a_{<i})$---is
\emph{independent of $T$}, so
\begin{multline*}
\Pr[q_i = T \mid Z = z]
= \sum_t \Pr[q_i = t]\,\Pr[T = t \mid Z = z] \\
\le \max_t \Pr[T = t \mid Z = z] = 2^{-k_z},
\end{multline*}
giving $\Pr[\exists i : q_i = T \mid z] \le q\,2^{-k_z}$. Taking the expectation
over $z$ and using $\mathbb{E}_z[2^{-k_z}] = \mathsf{Guess}(T \mid Z) = 2^{-k}$
from Definition \ref{def:k} yields $q \cdot 2^{-k}$. An adversary that outputs
an unqueried value wins with probability at most $2^{-k_z} \le q\,2^{-k_z}$,
which the same bound absorbs. The term $q \cdot 2^{-\lambda}$ charges the
probability that a query hits the uniform value $\beta(T)$ by chance.

\emph{(L).} Immediate from Definition \ref{def:k} after conditioning on
$\beta(T)$; no hypothesis on $\beta$ is used.
\end{proof}

\begin{remark}[The hypothesis in (H) is necessary, and our two instances straddle it]
\label{rem:betanecessary}
Without the hypothesis, part (H) is not merely unproven but \emph{false}: take
$\beta = \mathrm{id}$, or $\beta = \mathsf{encode}$ as in
Proposition \ref{prop:semantic}. The adversary outputs $\beta(T)$ and succeeds
with probability $1$, independently of $q$ and $k$. Our own two public values
therefore sit on opposite sides of the hypothesis---$b =
H(\mathsf{dom}\|n\|T)$ satisfies it, $\mathsf{encode}(T)$ does not---and this is
what makes the pair a theorem about \emph{our design} rather than a citation of
a standard bound. An earlier draft of this work stated (H) without its
hypothesis; the counterexample above is why the hypothesis is now explicit.
\end{remark}

\begin{corollary}[Re-execution, not enumeration]
\label{cor:reexec}
Under Assumption \ref{as:Z}, when the decoding is greedy or low-temperature,
$k \approx 0$ and the attack of Theorem \ref{thm:t1} is realised by
\emph{re-running the model}, at the cost of one inference---independent of
$|T|$ and requiring \emph{no} evaluation of $\beta$. Enumeration over the string
space is \emph{not} the relevant attack.
\end{corollary}

\fixme{原 docs/sp1-zk-audit.md §4 写的是"熵(~10^2--10^3 bit)远低于 256 bit，
因此枚举可行"——这是非因果推论（2^100 与 2^256 一样枚举不了；真正的攻击是\textbf{重执行}，
对手的已知信息是 (prompt, 模型, 解码参数)）。
2026-10-12 已全库改正，docs/tdsc-fixlist.md §G 的 G1--G6 逐条结案。}

\begin{remark}[A design observation, not a separation]
\label{rem:designated}
Our verification interface is \emph{designated}: binding is checked against a
value $b = H(\mathsf{dom}\|n\|T)$ in which a verifier-issued nonce $n$ appears.
It is tempting to promote this into a separation (``public verification is
impossible when $k$ is small''). \textbf{We claim no such theorem, and
Corollary \ref{cor:reexec} is why.} When $k \approx 0$, hiding fails for
\emph{every} scheme, designated or not, because the adversary re-runs the
model; when $k$ is large, a zero-knowledge proof of ``$\exists\, T : b =
H(\mathsf{dom}\|n\|T) \wedge J(\pi, T) = 1$'' is simultaneously publicly
verifiable and hiding. The dichotomy is thus governed by $k$ alone and carries
no independent content. What does survive---and what
\S\ref{sec:wp4-tension}
develops---is that \emph{third-party auditing of delivery} forces disclosure of
$T$, which is an informational remark rather than a theorem. We record the
demotion here because an earlier draft asserted the separation as a theorem.
\end{remark}

\need{\textbf{This is retrospective justification, not a novel construction.}
The system publishes $b = H(\mathsf{dom}\|n\|T)$ with a verifier-issued nonce
$n$, and Remark \ref{rem:designated} explains the shape of that choice. The
security model does not pin down whether $n$ is public; making it public
forfeits hiding, and making it verifier-held is what makes verification
designated.}

\need{\textbf{实验内容}：实测 $k$。对一组真实 prompt、同一模型多次采样，
估 $\mathsf{Guess}(T\mid Z)$（Definition \ref{def:k}），代入
Theorem \ref{thm:t1}(H) 得阈值。
据先例扫查，\emph{没有人测过这个量}——它本身是一块贡献。}

\begin{proposition}[The lossy instance of Theorem \ref{thm:t1}(L)]
\label{prop:semantic}
For the semantic rule class, the public values contain $\beta(T) =
\mathsf{encode}(T)$, a content-preserving feature vector of the response. This
$\beta$ is deterministic and injective on the features but not on strings, so
it violates the hypothesis of Theorem \ref{thm:t1}(H) and instantiates part
(L): the binding force is exactly
$\mathsf{Guess}(T \mid Z, \mathsf{encode}(T))$, i.e.\ the guessing probability
\emph{within the preimage class} of the published vector. Here $k$ need not be
argued from perplexity at all: the adversary's search space is that preimage,
which is dramatically smaller than the string space. The relevant asymmetry is
not ``low entropy'' but a \emph{lossy public encoding}.
\end{proposition}
\begin{proof}
Immediate. The measurement is direct: publish $\mathsf{encode}(T)$ for real
responses and measure how much of $T$ is recoverable.
\end{proof}

\need{This proposition is the paper's cheapest experiment and its strongest
anchor for Theorem \ref{thm:t1}---it is measurable in a day, and it makes the
bound concrete rather than cited. Source: \texttt{docs/design-semantic-rules.md}
lines 195--201, where the repository itself already names it as a second
instance.}

\subsection{Quantitative request binding}
\label{sec:t3}
% 决策 D2（第二版，2026-10-12）：原「闭合三分法」定理降级为 Remark（分类），
% 取代它的是下面的定量界 T3′。理由见 docs/tdsc-plan.md §1.6(b)：
% 三分法要么假（理想功能模型下能构造出第四条路），要么琐碎（把三条路写成定义）。

\begin{theorem}[Selective-presentation advantage]
\label{thm:t3}
Fix a policy $\pi$ and a request distribution. A selective-presentation
adversary $\mathcal{B}$ may, \emph{before} seeing the receiver's challenge,
compute a set $S$ of accepted transcripts. The receiver then contributes $f$
fresh, uniformly random bits $r$ that enter the proven statement. Then
\[
\Pr\!\left[\mathcal{B}\ \text{is accepted on } r\right] \;\le\; |S| \cdot 2^{-f}.
\]
If instead the receiver contributes no fresh randomness ($f = 0$; e.g.\ an
empty nonce), the advantage is $1$, and \emph{does not decay in any parameter
of the system}.
\end{theorem}
\begin{proof}
\need{主论证：$r$ 进入 $b = H(\text{dom} \| r \| T)$，故 $S$ 里任何一个
预先算好的证书，只有在 $r$ 恰好落进它已覆盖的挑战集合时才被接受；
$f$ 个新鲜均匀比特给出 $2^{-f}$ 的对撞上界，再对 $S$ 取并界。
$f = 0$ 时 $S$ 中每个元素都是永久有效的重放。}

\need{\textbf{代码依据}：\texttt{circuits/types/src/lib.rs}:950--995 已写明
“nonce 为空是合法的……不提供任何重放防护”。定理的 $f = 0$ 分支不是假想，
是本系统当前 API 的\emph{合法输入}。}
\end{proof}

\need{这条定理与 Theorem \ref{thm:t1} \textbf{同构}，而正是这一点使它值得写：
两条说的是同一件事——\emph{一个证明只能绑住提交进去的、对手无法预先选择的
那部分比特所能支撑的量}。Theorem \ref{thm:t1} 管响应侧
（$k = -\log_2 \mathsf{Guess}(T \mid Z)$，见 Definition \ref{def:k}），Theorem \ref{thm:t3} 管请求侧（新鲜度 $f$）。
这条统一主线写在 \S\ref{sec:intro} 的贡献列表里，是全文的理论骨架。}

\begin{remark}[A taxonomy, not a trichotomy]
\label{rem:taxonomy}
Every policy-proof deployment we surveyed closes the divergent-delivery gap of
Proposition \ref{thm:gap} in one of three ways:
(a) root trust in hardware (attested enclaves);
(b) have the \emph{end user} contribute a secret to the proven statement (the
user issues the nonce); or
(c) gate delivery on successful verification (\emph{prove-before-deliver}).
\textbf{We offer this as a classification, not as an exhaustiveness theorem}:
we neither prove that a fourth family cannot exist, nor that the three are
disjoint. zk-attach \cite{zkattach} walks (a); NiyamAI and ZKAP walk (c); PoP
walks (b).
\end{remark}

\need{\textbf{Posture, stated plainly.} Our system already carries
\texttt{response\_binding}, but $n$ is issued by the \emph{verifier}. Moving it
to the \emph{end user}---making $f > 0$ the receiver's contribution---closes
the divergent-delivery branch of Proposition \ref{thm:gap}. This is a small
protocol amendment, not a new system, and the paper must say exactly that. What
Theorem \ref{thm:t3} buys is the \emph{price}: with $f = 0$ there is no closure
at any cost.}

% =====================================================================
\section{Policy-Private Compliance}
\label{sec:private-policy}
% 决策 D3：这是"透明模式 ZK 冗余"这个批评的唯一硬解。
% 2026-10-12 读代码后的定性：这不是"造一个新证明系统"，是一次\emph{披露策略}
% 的重构——策略文本在架构上已经是证人。见 docs/tdsc-plan.md WP4.1。

\need{动机：对确定性规则 + 公开响应，签名哈希链就够了。唯一无法用签名
替代的场景是策略本身要保密。}

\subsection{The observation that reframes this work package}
\label{sec:wp4-finding}

\texttt{spec\_canonical} is already a \emph{guest input}---a witness---not a
public value: only
$\mathsf{policy\_hash} = \mathsf{SHA256}(\texttt{spec\_canonical})$
enters the public values (\texttt{circuits/types/src/lib.rs}:1--16, 415--416,
918, 1003--1004). So ``do not reveal the policy text'' is
\emph{architecturally already true}; what leaks is everything \emph{derived}
from the policy. WP4 is accordingly a \emph{disclosure-policy refactor}, not a
new proof system, and must be written that way---as a systematic account of
what a certificate leaks about the policy it certifies.

\begin{table}[t]
\centering
\caption{What a policy-private certificate leaks today, and what closing each
leak costs.}
\label{tab:wp4}
\begin{tabular}{@{}L{0.04\linewidth}L{0.46\linewidth}L{0.36\linewidth}@{}}
\toprule
\# & Leak & Difficulty / dependency \\
\midrule
1 & The policy text itself & \textbf{Already closed} (it is a witness). \\
2 & Witness hiding in the proof system & Medium; \emph{inherited from WP1}:
    only \texttt{groth16}/\texttt{plonk} is witness-hiding, per
    \texttt{docs/\allowbreak sp1-zk-audit.md}. \\
3 & \texttt{PrivateViolation}\allowbreak\{\texttt{rule},\allowbreak\texttt{kind}\} names the rule that
    fired & Medium; needs a coarsened violation class plus a commitment to it. \\
4 & Evidence opening reveals the rule
    (\texttt{policydsl/\allowbreak privacy/\allowbreak commit.py}:319) & Medium; disclose the
    \emph{existence} of a violation, not its identity. \\
5 & The \texttt{delegated} flag exposes policy structure (which rules were
    routed to the ezkl path) & Medium; easiest to overlook, cheapest to fix. \\
6 & Semantic rules publish $\beta(T)=\mathsf{encode}(T)$
    (Proposition \ref{prop:semantic}) & Medium--high; either bound the
    preimage, or disable this rule class in policy-private mode. \\
\bottomrule
\end{tabular}
\end{table}

\need{Honest structure of the table, and it should be stated as such:
sub-problem 1 is a one-paragraph re-reading of the existing code, \emph{not} a
contribution; 2 is inherited wholesale from WP1; 3--5 are engineering with one
real formal core (disclosure granularity); 6 is where the work is---and it
collides with Theorem \ref{thm:t1}, because one cannot hide the policy while
publishing a lossy encoding of the response that the semantic rules require.}

\subsection{The tension worth writing up: policy privacy vs.\ accountability}
\label{sec:wp4-tension}

Lead the section with the tension, not with the cryptography: it determines
whether the mode means anything at all.

\paragraph{The prerequisite, stated as a prerequisite.} Let
$H = \mathsf{SHA256}(\pi)$. On its own $H$ is \emph{inert}. A prover can always
emit a self-consistent certificate for any policy $\pi'$ it prefers, simply by
placing $H' = \mathsf{SHA256}(\pi')$ in that certificate. \emph{Secrecy of
$\pi$ is what creates the gap}, and no cryptographic mechanism can close it:
closing it would require showing $\pi$ to the verifier, which is precisely what
the mode exists to avoid. What closes it is an \emph{issuing act}, so we state
it as an explicit prerequisite of the mode:

\begin{description}
  \item[(P-PP)] An \emph{issuing authority} signs $\pi$ and registers $H$ in a
  record readable by third parties. A policy is \emph{accountable} iff its
  \texttt{policy\_hash} can be traced to such a record.
\end{description}

Without (P-PP) the mode is not merely weaker---the verifier's conclusion is
empty. This is the same assumption as ``honest policy author'' in
\texttt{docs/security-model.md}~\S0.2; what is new here is scoping it to the
verification interface, which had not been done, and which a referee will
otherwise use to dismiss the whole mode in one sentence.

\paragraph{Three parties.} ``Three'' means \emph{authority / operator / third
party}; an auditor is a named instance of the third party, not a fourth trust
anchor.

\begin{table}[t]
\centering
\caption{The three-party model of policy-private mode.}
\label{tab:wp4-parties}
\setlength{\tabcolsep}{3pt}
\begin{tabular}{@{}L{0.15\columnwidth}L{0.13\columnwidth}L{0.64\columnwidth}@{}}
\toprule
Party & Honest? & Holds; wants \\
\midrule
Authority $R$ & \textbf{Yes}, by (P-PP) & Holds $\pi$ and the issuing record;
wants $\pi$ actually enforced, and wants to be able to check that it was. \\
Operator $O$ & \textbf{No} (the adversary) & Holds $\pi$, $T$, the proving key;
wants $\pi$ hidden from third parties while still being able to certify. \\
Third party $V$ & Yes & Holds only public values---\texttt{policy\_\allowbreak hash},
\texttt{circuit\_\allowbreak hash}, \texttt{response\_\allowbreak commitment}, \texttt{passed}, the
proof; wants to confirm that \emph{some issued} policy passed, and must
\emph{not} learn $\pi$. \\
\emph{(Auditor $V^{*}$)} & Yes & As $V$, plus $\pi$ once $O$ discloses it;
wants ex-post accountability. \\
\bottomrule
\end{tabular}
\end{table}

Two boundaries must be stated with the model. First, secrecy of $\pi$ comes
from a \emph{data-flow fact}---$\pi$ never enters the public values---not from
zero-knowledge; any rule that writes a fragment of $\pi$ into the public values
(the leaks in Table~\ref{tab:wp4}, items~3 and~5) breaks it on that fragment.
Second, accountability for $V^{*}$ is \emph{voluntary}: $V^{*}$ can check only
when $O$ chooses to disclose $\pi$. The mode offers no forced disclosure; that
would need the release-gating family discussed in
\texttt{docs/tdsc-plan.md}~\S1.6(b), which we do not claim.

\paragraph{Verifier capability table.} This table is the contribution; the
cryptography is downstream of it, and the table survives even if the
cryptography only partially lands. \emph{It must appear in the paper as
written}, because it is what stops \texttt{policy\_hash} from being read as
``compliance was proved''.

\begin{table}[t]
\centering
\caption{What each verifier may and may not conclude in policy-private mode.}
\label{tab:wp4-capability}
\footnotesize
\setlength{\tabcolsep}{3pt}
\begin{tabular}{@{}L{0.14\columnwidth}L{0.31\columnwidth}L{0.30\columnwidth}L{0.17\columnwidth}@{}}
\toprule
Verifier & \textbf{May} conclude & \textbf{May not} conclude & Rests on \\
\midrule
Third party $V$ & (i) the commitment binds $T$
(Definition~\ref{def:k}); (ii) \emph{some} policy with hash $H$ passed;
(iii) $H$ appears in the issuing record & (iv) the \emph{content} of $\pi$;
(v) how many rules $\pi$ has, or of which classes; (vi) ``the operator
enforces the regulation \emph{I} care about''---only ``some issued
regulation'' & (P-PP); A2; the $k$ of T1 \\
End user (also receiver) & As $V$, and additionally: the proved $T$ is the
string it received, since its nonce is bound (T3$'$, requires $f>0$) & Only
the string it received; \emph{not} what was delivered to another downstream
agent & $f>0$ \\
Authority $R$ & Everything; it holds $\pi$ and can recompute $H$ & --- & --- \\
Auditor $V^{*}$ (after disclosure) & (vii) $\mathsf{SHA256}(\pi) =
\texttt{policy\_\allowbreak hash}$, i.e.\ what ran was the issued $\pi$;
(viii) the disclosed violation record is consistent & (ix) the undisclosed
part of $T$ & voluntary disclosure \\
On-chain verifier & (i)--(iii) only; it hashes and verifies signatures, and
does not interpret policy semantics & (iv)--(ix) & (P-PP) \\
\bottomrule
\end{tabular}
\end{table}

The $V$ row is the whole reason this subsection exists: $V$ does not obtain
``compliance'', it obtains ``some issued policy was satisfied''. Conflating the
two is the sentence a referee would use to kill the mode.

\need{Remaining work on this subsection: (a) confirm the (P-PP) framing with
the collaborator, and choose its instantiation (an on-chain registry of issued
policy hashes vs.\ a signed manifest); (b) add the prior-art citations below;
(c) decide whether the capability table merges with Table~\ref{tab:wp4} in the
submission, or stays separate (currently separate, on the grounds that one
table is about leakage and the other about verifier conclusions).}

\subsection{Expected results, stated in advance}
\label{sec:wp4-expect}

\need{Calibrated expectations, so that a partial outcome is a publishable
result rather than a failure:}
\begin{itemize}
  \item Sub-problems 1, 2, 6: high confidence ($\gtrsim 80\%$). 1 is already
        true of the code; 2 follows from WP1; 6 is a design restriction.
  \item Sub-problems 3, 4, 5: likely ($\sim 70\%$).
  \item ``\emph{Zero}-knowledge violation evidence'' in the strong sense:
        only $50$--$60\%$. The boolean ``some violation occurred'' is easy; the
        hard part is \emph{granularity}---with $3$--$8$ rules, a coarsened
        violation class can itself identify the rule. \textbf{The likely best
        result is a granularity/actionability/leakage tradeoff curve, not full
        hiding}---and that curve is a contribution in its own right.
  \item \S\ref{sec:wp4-tension}, the tension and the capability table:
        $\sim 90\%$. This is precisely why the section leads with it.
\end{itemize}

\need{Prior-art discipline, and this is the highest-risk place in the paper ---
the collaborator's own field. \emph{Policy hiding is not a new primitive}:
hidden credentials, secret handshake, and XACML privacy in access control, and
attribute-hiding / fuzzy predicate encryption, are all the same family. Cite
that line explicitly and do not claim primacy. What is new is narrow: (a) the
\emph{application} to LLM output compliance together with selective disclosure
of violation type and location, and (b) the accountability tension of
\S\ref{sec:wp4-tension}.}

% =====================================================================
\section{Composition: The Constructive Half of the End-User-Secret Family}
\label{sec:composition}
% 决策 D4：双载体（ezkl + 真 ONNX 主臂；DeepProve 对照臂）。

\fixme{关键更正：组合本身不闭合缺口。只有当模型证明绑定到\emph{用户提交的
输入} $x$ 时，$M(x) = T \wedge \pi(T)$ 才真的闭合"送错字节"——而那恰好是
Theorem \ref{thm:t3} 要求的 $f > 0$：用户提交的 $x$ 就是那 $f$ 个新鲜比特。}

\need{组合契约的定义（模型证明须暴露什么，策略证明才能绑上）；
真模型规模（$\ge 10^6$ 参数）；成本数据。}

% =====================================================================
\section{Availability: The Unanchored-Pattern Cost Attack}
\label{sec:dos}
% WP6。

\begin{theorem}[Quadratic cost inflation]
\need{对含未锚定模式的正则，控制响应的对手可强制 $\Omega(n^2)$ 证明工作量。}
\end{theorem}

\need{缓解：Pike VM（已有）+ 响应长度承诺 + 分块匹配。
zk-regex 有 DFA 限制但没有 zkVM 场景下的 DoS 分析——这是新的。}

% =====================================================================
\section{Coverage: Can a Provider Simply Withhold Certificates?}
\label{sec:coverage}
% WP5。诚实预期：无可信硬件时可能无解 —— 那就写成 open problem。

% =====================================================================
\section{Evaluation}
\label{sec:eval}

\need{矩阵：规则 \{1,2,3,4,6,8\} $\times$ 长度 \{200, 2k, 10k, 50k\}
$\times$ 模式 \{transparent, private, private-policy\}
$\times$ 证明档 \{core, groth16\}。
骨架见 bench/bench_scale.py。}

\subsection{Why zero-knowledge? A decision map}
% WP9：非 ZK 基线 + 决策图。

\need{签名 Merkle 日志对照臂（brivora/VAOL 风格）。
结论：ZK 只在 \{私有响应, 私有策略\} 两个角上有正收益。
把"透明模式 ZK 冗余"从批评变成有刻度的设计结论。}

% =====================================================================
\section{Limitations}
\label{sec:limitations}

\need{必须写进去的：下列五条。每条都要在正文里展开，不是罗列。}
\begin{itemize}
  \item Content privacy here is \emph{commitment} privacy, not
        irrecoverability of content (Corollary \ref{cor:reexec}).
  \item The zero-knowledge property of the default proof mode:
        \fixme{SP1's published security model states explicitly that
        individual STARK proofs do not currently satisfy the
        zero-knowledge property. If groth16 does not become the delivery
        mode, the word ``zero-knowledge'' must leave the title.}
  \item Binding the composition's input requires the user to be online (or to
        hold a key). This is the price of Remark \ref{rem:taxonomy}(b), and
        the concrete form of $f > 0$ in Theorem \ref{thm:t3}.
  \item The coverage problem (\S\ref{sec:coverage}) may be unsolvable without
        trusted hardware.
  \item Three-party collusion is undetectable at the signature layer
        (\texttt{docs/security-model.md}).
\end{itemize}

% =====================================================================
\section{Conclusion}

% =====================================================================
\appendix
\section{Proofs}
\need{\textbf{状态（2026-10-12 自检后）}：
Theorem \ref{thm:t1}(H) 与 Theorem \ref{thm:t3} 的证明\textbf{已写在正文}，
此处不再重复。三分法已降级为分类 Remark \ref{rem:taxonomy}，无需证明。
Theorem \ref{thm:t1}(L) 无闭式，只给测量口径（\S\ref{sec:eval}）。
原 T2「分离定理」已\textbf{停用}（见 Remark \ref{rem:designated}）——
本节因此只剩一条待补事项，见下。}

\need{\textbf{待补}：把 Theorem \ref{thm:t1}(H) 的随机预言机模型换成
标准模型（或明确声明本文只在 ROM 下主张）。当前证明依赖 $H$ 是新鲜随机函数
这一条；实践中的 SHA-256 只给计算意义下的近似，故正文措辞必须写成
\emph{heuristic}，不得写成信息论界。}

% =====================================================================
\begin{thebibliography}{99}

\bibitem{hollowllm}
C.~Gong, B.~Liu, and M.~Li, ``Hollow-LLM Attack: Computationally Trivial
Weights in Zero-Knowledge Verification of LLM Inference,'' in \emph{IEEE S\&P},
2026, pp.~4339--4355, DOI 10.1109/SP63933.2026.00258.

\bibitem{zkattach}
\need{zk-attach, arXiv:2608.30083 --- 打开原文核实作者后再引}

\bibitem{zkgen}
J.~Lauinger, J.~Ernstberger, and S.~Steinhorst, ``zkGen: Policy-to-Circuit
Transpiler,'' in \emph{IEEE ICBC}, 2024, pp.~620--624,
DOI 10.1109/ICBC59979.2024.10634440.

\bibitem{pco}
\need{PCO, IACR ePrint 2026/994 --- 打开原文核实作者（Dolev/Yung 等）后再引}

\bibitem{zkagent}
L.~Wang, H.~Lou, C.~Li, Y.~Yu, and Y.~Hu, ``zkAgent: Verifiable LLM Agent
Execution via One-Shot Transcript Proofs,'' IACR ePrint 2026/199.

\bibitem{vdr}
\need{Verifiable Document Redacting, Springer, DOI 10.1007/978-3-319-66402-6_20
--- 核实作者与年份。\textbf{绝不能再引 IACR ePrint 2024/813}（那是
Kiraz/Larraia/Vaughan, "How to Redact the Bitcoin Backbone Protocol"）。
见 docs/tdsc-fixlist.md \S A。}

\end{thebibliography}

\end{document}
"   % 两次，第二遍收交叉引用
%  \SKELQUIET 让所有中文标记宏吞掉参数；不带它则需要能排版中文的引擎
%  （xelatex + CJK 字体），本骨架未配置。
% =====================================================================
\documentclass[journal]{IEEEtran}

\usepackage{amsmath,amssymb,amsthm}
\usepackage{booktabs}
\usepackage{array}
\usepackage{graphicx}
\usepackage{xcolor}
\usepackage{hyperref}

% 窄栏里用 p{} 会因两端对齐产生大片字距（编译时刷一堆 Underfull \hbox）。
% 表格正文列一律用 L{}：左对齐、允许自然参差。
\newcolumntype{L}[1]{>{\raggedright\arraybackslash}p{#1}}
\newcommand{\tabfont}{\footnotesize}

\newtheorem{theorem}{Theorem}
\newtheorem{lemma}[theorem]{Lemma}
\newtheorem{proposition}[theorem]{Proposition}
\newtheorem{corollary}[theorem]{Corollary}
\newtheorem{definition}[theorem]{Definition}
\newtheorem{assumption}[theorem]{Assumption}
\newtheorem{remark}[theorem]{Remark}

% --- 工作标记：提交前全部清掉 ---
% 注意：宏名不能叫 \wp —— \wp 已是 LaTeX/amssymb 的 Weierstrass ℘，
% \newcommand{\wp} 会直接报 "Command \wp already defined" 并使后文级联报错。
\ifdefined\SKELQUIET
  \newcommand{\wpmark}[1]{}
  \newcommand{\need}[1]{}
  \newcommand{\fixme}[1]{}
\else
  \newcommand{\wpmark}[1]{\textcolor{red}{\textsuperscript{[WP#1]}}}
  \newcommand{\need}[1]{\textcolor{blue}{[[ NEED: #1 ]]}}
  \newcommand{\fixme}[1]{\textcolor{orange}{[[ FIX: #1 ]]}}
\fi

\begin{document}

% ---------------------------------------------------------------------
% 标题：决策 D5 —— 保留 "Zero-Knowledge" 的前提是 groth16 成为交付档 (WP1)。
%       若 WP1 未打通 groth16，必须删掉标题与摘要里的 zero-knowledge。
% ---------------------------------------------------------------------
\title{Proof-of-Policy: Zero-Knowledge Compliance Certificates\\ for LLM Agent Outputs}
% 备选标题（若 groth16 未打通）："Proof-of-Policy: Succinct Compliance
% Certificates for LLM Agent Outputs"

\author{\need{authors}}
\maketitle

% ---------------------------------------------------------------------
\begin{abstract}
% 决策 E2：内容隐私上界必须出现在摘要，不能埋进 §8.1。
We present \emph{Proof-of-Policy} (PoP), a system that proves an LLM agent's
response $T$ satisfies a publicly declared policy $\pi$, without a trusted
execution environment...
\need{一句完整的贡献陈述（对应 D1 的"分析 x 系统各半"）}

\smallskip
\noindent\textbf{Scope of privacy.} PoP provides \emph{policy}-zero-knowledge:
the certificate reveals neither the policy evaluation trace nor, in
policy-private mode, the policy itself. It does \emph{not} guarantee that $T$
is unrecoverable: as we show in \S\ref{sec:binding}, any publicly recomputable
binding to $T$ admits a re-execution attack whenever the adversary knows the
prompt and the model. The same bit-budget governs the request side: a
certificate is replayable without limit unless the receiver contributes fresh
randomness (Theorem \ref{thm:t3}).
\end{abstract}

\begin{IEEEkeywords}
zero-knowledge proofs, LLM agents, policy compliance, verifiable computation,
accountability
\end{IEEEkeywords}

% =====================================================================
\section{Introduction}
\label{sec:intro}

\need{Contribution list. 必须包含下面这条"不主张"声明：}

\noindent\textbf{What we do \emph{not} claim.}
We do not claim that policy proofs bind nothing: our own system already binds
the response via a verifier-issued nonce (Definition 1). We do not claim to be
the first to observe that a proof can fail to bind the artifact it purports to
describe: the \emph{effort gap} for ZK inference proofs is formalised by
Hollow-LLM \cite{hollowllm}, and the \emph{source-integrity gap} for predicate
proofs is named by zk-attach \cite{zkattach}. Our contribution is to
(i) show that the surviving gap is the \emph{divergent-delivery} one
(Proposition \ref{thm:gap}), (ii) give a quantitative bound on what a public
value reveals, together with the \emph{structural hypothesis} under which that
bound holds---and a demonstration that the hypothesis is necessary, since our
own two public values sit on opposite sides of it (Theorem \ref{thm:t1} and
Remark \ref{rem:betanecessary}), and (iii) give a quantitative price for
closing the gap (Theorem \ref{thm:t3}).

\smallskip
\noindent\textbf{A separation we deliberately do \emph{not} claim.} It is
tempting to promote our designated-verifier interface into a theorem
(``public verification is impossible when $k$ is small''). We give the reason
we do not in Remark \ref{rem:designated}: in the low-$k$ regime hiding fails
for every scheme, designated or not, and the dichotomy has no content of its
own. We state this explicitly because the negated claim is the more attractive
one to write.

\smallskip
\noindent\textbf{One spine, two sides.} Theorems \ref{thm:t1} and
\ref{thm:t3} state the same fact about two different inputs: \emph{a proof
binds only as much as the committed, adversary-non-selectable quantity has bits
for}. The response side is governed by entropy,
$k = -\log \mathsf{Guess}(T \mid Z)$ (Definition \ref{def:k}); the request side
by freshness, $f$. Neither can be bought with more computation.

\fixme{全文禁止出现"比 Proof-of-Guardrail 更强"或任何强弱对比。
Proof-of-Guardrail 与本文正交：它证"计算发生过"，本文证"性质成立"。}

% =====================================================================
\section{Background and Related Work}
\label{sec:related}

\need{五条线，每条都要"它占什么位 / 我们差在哪"两列。见 docs/tdsc-plan.md \S2。}

\begin{table}[t]
\centering
\caption{Positioning against the five prior-art lineages.}
\label{tab:related}
\setlength{\tabcolsep}{3pt}
\begin{tabular}{@{}L{0.26\columnwidth}L{0.42\columnwidth}L{0.26\columnwidth}@{}}
\toprule
Lineage & Representative work & Occupied slot \\
\midrule
Policy $\to$ circuit      & zkGen \cite{zkgen}, Circuitree, Prezta & DSL-to-circuit \\
Agent/LLM compliance      & NiyamAI, Zaguk, PCO \cite{pco}, WITNESS & private-mode kernel \\
Model computation         & DeepProve, zkGPT, zkAgent \cite{zkagent} & inference integrity \\
Attacks / negative results& Hollow-LLM \cite{hollowllm}, zk-attach \cite{zkattach} & naming the gap \\
Non-ZK industry baseline  & @brivora/verify, VAOL & signed audit logs \\
\bottomrule
\end{tabular}
\end{table}

\fixme{引文纪律：除 Hollow-LLM / zkGen / NiyamAI / zkAgent / ePrint 2024/813
五条外，其余引文必须逐条打开原文核实后再进 bibliography。
见 docs/tdsc-fixlist.md \S2 的注记。}

% =====================================================================
\section{System Overview}
\label{sec:system}
% 决策 D1 的"系统半边"：这一节压缩到 ~1.5 页，是 T3(b) 的载体，不是主角。

\need{DSL $\to$ ConstraintSpec $\to$ circuit；双模式；证书 + 锚定；
工具回执链；801 测试 / 15 skip 的工件陈述。}

% =====================================================================
\section{Binding of Policy Proofs}
\label{sec:binding}
% 决策 D1 的"分析半边"：这是论文的核心。

\subsection{The gap}

\begin{definition}[Policy-proof system]
A policy-proof system is a tuple $(\mathsf{Prove}, \mathsf{Ver}, \pi, \mathcal{M})$.
On challenge $n$ from a verifier, $\mathsf{Prove}(\pi, T, n)$ outputs a
certificate $c$ whose public values include
\begin{multline*}
b \;=\; H(\texttt{pop-bind-v1} \,\|\, \mathsf{u32be}(|n|) \,\|\, n \,\|\, T),\\
C = H(T).
\end{multline*}
$\mathsf{Ver}(c, n, T')$ accepts iff the policy relation holds \emph{and}
$b$ recomputed from the received $T'$ matches.
\end{definition}

\fixme{\textbf{H3/H4}：早先写的 ``the relation contains no term referring to the
delivered bytes'' 与仓库事实冲突——公开值里\emph{已有} $b$，验证时由送达的
$T'$ 现场重算（\texttt{circuits/types/src/lib.rs:444}、
\texttt{docs/security-model.md:232}）。命题已按上面的定义改写。见
\texttt{docs/tdsc-fixlist.md} \S H。}

\begin{proposition}[Divergent delivery]
\label{thm:gap}
Definition 1 binds $T'$ \emph{as received by the verifying party}. If the
provider delivers a different string $T'' \ne T'$ to the end user, verification
still succeeds on $T'$. Neither hardware nor a proof can prevent this; only the
end user holding $n$ and recomputing $b$ can.
\end{proposition}
\begin{proof}
$b$ is a function of $(n,T)$ only; the delivery channel to the end user carries
no authenticated term. One line.
\end{proof}

\need{This is the \emph{residual} gap, not ``proofs bind nothing''. Prior work
must be cited here (Hollow-LLM \cite{hollowllm}, zk-attach \cite{zkattach}), and
we state plainly that the naive claim would be false against our own system.}

\subsection{A quantitative binding--hiding tradeoff}
\label{sec:tradeoff}

% ---------------------------------------------------------------------
% T1 : 关键是把"对手已知信息 Z"定对。
%      docs/tdsc-plan.md §1.4 已修正原 §4 的非因果推理。
% ---------------------------------------------------------------------
\begin{definition}[Guessing entropy of a response]
\label{def:k}
For a response distribution $\mathcal{D}$ and adversary side information $Z$,
let
\[
\mathsf{Guess}(T \mid Z) \;\triangleq\; \max_{\mathcal{A}} \Pr[\mathcal{A}(Z) = T],
\]
which equals $\mathbb{E}_z[\max_t \Pr[T = t \mid Z = z]]$, and
\[
k \;\triangleq\; -\log_2 \mathsf{Guess}(T \mid Z),
\]
the maximum ranging over all time-unbounded adversaries.
\end{definition}

\begin{remark}[The definition is operational, and it is the average]
\label{rem:kmeas}
$k$ is \emph{measured}, not assumed: sampling one model on one prompt many
times estimates $\mathsf{Guess}(T \mid Z)$ directly
(\S\ref{sec:eval}), so Theorem \ref{thm:t1} is instantiated with a
number we obtain rather than a number we posit. The definition must average
over $z$ rather than take the pointwise worst case
$k_z = -\log \max_t \Pr[T = t \mid Z = z]$: the proof takes an expectation over
$z$, and $\mathbb{E}_z[2^{-k_z}] = 2^{-k}$ holds only for the average form.
\end{remark}

\begin{assumption}[Adversary side information]
\label{as:Z}
The adversary knows the prompt, the model (possibly its weights), and the
decoding parameters. \emph{This is the assumption that makes the theorem
non-trivial}: under greedy decoding, $k \approx 0$.
\end{assumption}

\begin{theorem}[Binding, and the hypothesis it needs]
\label{thm:t1}
Let $\beta : \mathcal{M} \to \{0,1\}^\lambda$ be public, deterministic and
efficiently computable, and let the certificate publish $\beta(T)$. Let
$\mathcal{A}$ be any adversary, unbounded in time, that is given $Z$ and
evaluates $\beta$ at most $q$ times.
\begin{enumerate}
\item[(H)] \emph{(Randomised public value.)} Suppose
  $\beta(T) = H(\mathsf{dom} \,\|\, \mathsf{u32be}(|n|) \,\|\, n \,\|\, T)$ for a
  random oracle $H$ and a nonce $n$ drawn uniformly and \emph{independently of
  $T$}. Then
  \[
  \Pr\!\left[\mathcal{A}^\beta\!\left(Z, \beta(T)\right) = T\right]
  \;\le\; q \cdot 2^{-k} \;+\; q \cdot 2^{-\lambda}.
  \]
  Consequently hiding with advantage $\varepsilon$ requires
  $k \ge \log q + \log(1/\varepsilon)$.
\item[(L)] \emph{(Lossy public value.)} For an arbitrary deterministic public
  $\beta$ the correct bound is over the \emph{preimage class}:
  \begin{multline*}
  \Pr\!\left[\mathcal{A}^\beta\!\left(Z, \beta(T)\right) = T\right]
  \;\le\; \mathsf{Guess}\!\left(T \mid Z, \beta(T)\right) \\
  \triangleq\; \mathbb{E}_y\!\left[\max_t \Pr[T = t \mid Z, \beta(T) = y]\right].
  \end{multline*}
\end{enumerate}
\end{theorem}

\begin{proof}
\emph{(H).} Fix $z$ and let the adversary's query set be
$Q = \{q_1, \dots, q_q\}$. Decompose by first hit,
\[
\Pr[\exists i : q_i = T] \;=\; \sum_i \Pr\!\left[q_i = T \wedge \neg E_{<i}\right],
\]
where $E_{<i} \triangleq \{\exists j < i : q_j = T\}$.
Conditioned on $\neg E_{<i}$, every oracle answer $a_1,\dots,a_{i-1}$ already
revealed is uniform and independent of $T$ (no queried point equals $T$, and $H$
is a fresh random function), and $\beta(T)$ itself is uniform and independent of
$T$. Hence $q_i$---a deterministic function of $(z, \beta(T), a_{<i})$---is
\emph{independent of $T$}, so
\begin{multline*}
\Pr[q_i = T \mid Z = z]
= \sum_t \Pr[q_i = t]\,\Pr[T = t \mid Z = z] \\
\le \max_t \Pr[T = t \mid Z = z] = 2^{-k_z},
\end{multline*}
giving $\Pr[\exists i : q_i = T \mid z] \le q\,2^{-k_z}$. Taking the expectation
over $z$ and using $\mathbb{E}_z[2^{-k_z}] = \mathsf{Guess}(T \mid Z) = 2^{-k}$
from Definition \ref{def:k} yields $q \cdot 2^{-k}$. An adversary that outputs
an unqueried value wins with probability at most $2^{-k_z} \le q\,2^{-k_z}$,
which the same bound absorbs. The term $q \cdot 2^{-\lambda}$ charges the
probability that a query hits the uniform value $\beta(T)$ by chance.

\emph{(L).} Immediate from Definition \ref{def:k} after conditioning on
$\beta(T)$; no hypothesis on $\beta$ is used.
\end{proof}

\begin{remark}[The hypothesis in (H) is necessary, and our two instances straddle it]
\label{rem:betanecessary}
Without the hypothesis, part (H) is not merely unproven but \emph{false}: take
$\beta = \mathrm{id}$, or $\beta = \mathsf{encode}$ as in
Proposition \ref{prop:semantic}. The adversary outputs $\beta(T)$ and succeeds
with probability $1$, independently of $q$ and $k$. Our own two public values
therefore sit on opposite sides of the hypothesis---$b =
H(\mathsf{dom}\|n\|T)$ satisfies it, $\mathsf{encode}(T)$ does not---and this is
what makes the pair a theorem about \emph{our design} rather than a citation of
a standard bound. An earlier draft of this work stated (H) without its
hypothesis; the counterexample above is why the hypothesis is now explicit.
\end{remark}

\begin{corollary}[Re-execution, not enumeration]
\label{cor:reexec}
Under Assumption \ref{as:Z}, when the decoding is greedy or low-temperature,
$k \approx 0$ and the attack of Theorem \ref{thm:t1} is realised by
\emph{re-running the model}, at the cost of one inference---independent of
$|T|$ and requiring \emph{no} evaluation of $\beta$. Enumeration over the string
space is \emph{not} the relevant attack.
\end{corollary}

\fixme{原 docs/sp1-zk-audit.md §4 写的是"熵(~10^2--10^3 bit)远低于 256 bit，
因此枚举可行"——这是非因果推论（2^100 与 2^256 一样枚举不了；真正的攻击是\textbf{重执行}，
对手的已知信息是 (prompt, 模型, 解码参数)）。
2026-10-12 已全库改正，docs/tdsc-fixlist.md §G 的 G1--G6 逐条结案。}

\begin{remark}[A design observation, not a separation]
\label{rem:designated}
Our verification interface is \emph{designated}: binding is checked against a
value $b = H(\mathsf{dom}\|n\|T)$ in which a verifier-issued nonce $n$ appears.
It is tempting to promote this into a separation (``public verification is
impossible when $k$ is small''). \textbf{We claim no such theorem, and
Corollary \ref{cor:reexec} is why.} When $k \approx 0$, hiding fails for
\emph{every} scheme, designated or not, because the adversary re-runs the
model; when $k$ is large, a zero-knowledge proof of ``$\exists\, T : b =
H(\mathsf{dom}\|n\|T) \wedge J(\pi, T) = 1$'' is simultaneously publicly
verifiable and hiding. The dichotomy is thus governed by $k$ alone and carries
no independent content. What does survive---and what
\S\ref{sec:wp4-tension}
develops---is that \emph{third-party auditing of delivery} forces disclosure of
$T$, which is an informational remark rather than a theorem. We record the
demotion here because an earlier draft asserted the separation as a theorem.
\end{remark}

\need{\textbf{This is retrospective justification, not a novel construction.}
The system publishes $b = H(\mathsf{dom}\|n\|T)$ with a verifier-issued nonce
$n$, and Remark \ref{rem:designated} explains the shape of that choice. The
security model does not pin down whether $n$ is public; making it public
forfeits hiding, and making it verifier-held is what makes verification
designated.}

\need{\textbf{实验内容}：实测 $k$。对一组真实 prompt、同一模型多次采样，
估 $\mathsf{Guess}(T\mid Z)$（Definition \ref{def:k}），代入
Theorem \ref{thm:t1}(H) 得阈值。
据先例扫查，\emph{没有人测过这个量}——它本身是一块贡献。}

\begin{proposition}[The lossy instance of Theorem \ref{thm:t1}(L)]
\label{prop:semantic}
For the semantic rule class, the public values contain $\beta(T) =
\mathsf{encode}(T)$, a content-preserving feature vector of the response. This
$\beta$ is deterministic and injective on the features but not on strings, so
it violates the hypothesis of Theorem \ref{thm:t1}(H) and instantiates part
(L): the binding force is exactly
$\mathsf{Guess}(T \mid Z, \mathsf{encode}(T))$, i.e.\ the guessing probability
\emph{within the preimage class} of the published vector. Here $k$ need not be
argued from perplexity at all: the adversary's search space is that preimage,
which is dramatically smaller than the string space. The relevant asymmetry is
not ``low entropy'' but a \emph{lossy public encoding}.
\end{proposition}
\begin{proof}
Immediate. The measurement is direct: publish $\mathsf{encode}(T)$ for real
responses and measure how much of $T$ is recoverable.
\end{proof}

\need{This proposition is the paper's cheapest experiment and its strongest
anchor for Theorem \ref{thm:t1}---it is measurable in a day, and it makes the
bound concrete rather than cited. Source: \texttt{docs/design-semantic-rules.md}
lines 195--201, where the repository itself already names it as a second
instance.}

\subsection{Quantitative request binding}
\label{sec:t3}
% 决策 D2（第二版，2026-10-12）：原「闭合三分法」定理降级为 Remark（分类），
% 取代它的是下面的定量界 T3′。理由见 docs/tdsc-plan.md §1.6(b)：
% 三分法要么假（理想功能模型下能构造出第四条路），要么琐碎（把三条路写成定义）。

\begin{theorem}[Selective-presentation advantage]
\label{thm:t3}
Fix a policy $\pi$ and a request distribution. A selective-presentation
adversary $\mathcal{B}$ may, \emph{before} seeing the receiver's challenge,
compute a set $S$ of accepted transcripts. The receiver then contributes $f$
fresh, uniformly random bits $r$ that enter the proven statement. Then
\[
\Pr\!\left[\mathcal{B}\ \text{is accepted on } r\right] \;\le\; |S| \cdot 2^{-f}.
\]
If instead the receiver contributes no fresh randomness ($f = 0$; e.g.\ an
empty nonce), the advantage is $1$, and \emph{does not decay in any parameter
of the system}.
\end{theorem}
\begin{proof}
\need{主论证：$r$ 进入 $b = H(\text{dom} \| r \| T)$，故 $S$ 里任何一个
预先算好的证书，只有在 $r$ 恰好落进它已覆盖的挑战集合时才被接受；
$f$ 个新鲜均匀比特给出 $2^{-f}$ 的对撞上界，再对 $S$ 取并界。
$f = 0$ 时 $S$ 中每个元素都是永久有效的重放。}

\need{\textbf{代码依据}：\texttt{circuits/types/src/lib.rs}:950--995 已写明
“nonce 为空是合法的……不提供任何重放防护”。定理的 $f = 0$ 分支不是假想，
是本系统当前 API 的\emph{合法输入}。}
\end{proof}

\need{这条定理与 Theorem \ref{thm:t1} \textbf{同构}，而正是这一点使它值得写：
两条说的是同一件事——\emph{一个证明只能绑住提交进去的、对手无法预先选择的
那部分比特所能支撑的量}。Theorem \ref{thm:t1} 管响应侧
（$k = -\log_2 \mathsf{Guess}(T \mid Z)$，见 Definition \ref{def:k}），Theorem \ref{thm:t3} 管请求侧（新鲜度 $f$）。
这条统一主线写在 \S\ref{sec:intro} 的贡献列表里，是全文的理论骨架。}

\begin{remark}[A taxonomy, not a trichotomy]
\label{rem:taxonomy}
Every policy-proof deployment we surveyed closes the divergent-delivery gap of
Proposition \ref{thm:gap} in one of three ways:
(a) root trust in hardware (attested enclaves);
(b) have the \emph{end user} contribute a secret to the proven statement (the
user issues the nonce); or
(c) gate delivery on successful verification (\emph{prove-before-deliver}).
\textbf{We offer this as a classification, not as an exhaustiveness theorem}:
we neither prove that a fourth family cannot exist, nor that the three are
disjoint. zk-attach \cite{zkattach} walks (a); NiyamAI and ZKAP walk (c); PoP
walks (b).
\end{remark}

\need{\textbf{Posture, stated plainly.} Our system already carries
\texttt{response\_binding}, but $n$ is issued by the \emph{verifier}. Moving it
to the \emph{end user}---making $f > 0$ the receiver's contribution---closes
the divergent-delivery branch of Proposition \ref{thm:gap}. This is a small
protocol amendment, not a new system, and the paper must say exactly that. What
Theorem \ref{thm:t3} buys is the \emph{price}: with $f = 0$ there is no closure
at any cost.}

% =====================================================================
\section{Policy-Private Compliance}
\label{sec:private-policy}
% 决策 D3：这是"透明模式 ZK 冗余"这个批评的唯一硬解。
% 2026-10-12 读代码后的定性：这不是"造一个新证明系统"，是一次\emph{披露策略}
% 的重构——策略文本在架构上已经是证人。见 docs/tdsc-plan.md WP4.1。

\need{动机：对确定性规则 + 公开响应，签名哈希链就够了。唯一无法用签名
替代的场景是策略本身要保密。}

\subsection{The observation that reframes this work package}
\label{sec:wp4-finding}

\texttt{spec\_canonical} is already a \emph{guest input}---a witness---not a
public value: only
$\mathsf{policy\_hash} = \mathsf{SHA256}(\texttt{spec\_canonical})$
enters the public values (\texttt{circuits/types/src/lib.rs}:1--16, 415--416,
918, 1003--1004). So ``do not reveal the policy text'' is
\emph{architecturally already true}; what leaks is everything \emph{derived}
from the policy. WP4 is accordingly a \emph{disclosure-policy refactor}, not a
new proof system, and must be written that way---as a systematic account of
what a certificate leaks about the policy it certifies.

\begin{table}[t]
\centering
\caption{What a policy-private certificate leaks today, and what closing each
leak costs.}
\label{tab:wp4}
\begin{tabular}{@{}L{0.04\linewidth}L{0.46\linewidth}L{0.36\linewidth}@{}}
\toprule
\# & Leak & Difficulty / dependency \\
\midrule
1 & The policy text itself & \textbf{Already closed} (it is a witness). \\
2 & Witness hiding in the proof system & Medium; \emph{inherited from WP1}:
    only \texttt{groth16}/\texttt{plonk} is witness-hiding, per
    \texttt{docs/\allowbreak sp1-zk-audit.md}. \\
3 & \texttt{PrivateViolation}\allowbreak\{\texttt{rule},\allowbreak\texttt{kind}\} names the rule that
    fired & Medium; needs a coarsened violation class plus a commitment to it. \\
4 & Evidence opening reveals the rule
    (\texttt{policydsl/\allowbreak privacy/\allowbreak commit.py}:319) & Medium; disclose the
    \emph{existence} of a violation, not its identity. \\
5 & The \texttt{delegated} flag exposes policy structure (which rules were
    routed to the ezkl path) & Medium; easiest to overlook, cheapest to fix. \\
6 & Semantic rules publish $\beta(T)=\mathsf{encode}(T)$
    (Proposition \ref{prop:semantic}) & Medium--high; either bound the
    preimage, or disable this rule class in policy-private mode. \\
\bottomrule
\end{tabular}
\end{table}

\need{Honest structure of the table, and it should be stated as such:
sub-problem 1 is a one-paragraph re-reading of the existing code, \emph{not} a
contribution; 2 is inherited wholesale from WP1; 3--5 are engineering with one
real formal core (disclosure granularity); 6 is where the work is---and it
collides with Theorem \ref{thm:t1}, because one cannot hide the policy while
publishing a lossy encoding of the response that the semantic rules require.}

\subsection{The tension worth writing up: policy privacy vs.\ accountability}
\label{sec:wp4-tension}

Lead the section with the tension, not with the cryptography: it determines
whether the mode means anything at all.

\paragraph{The prerequisite, stated as a prerequisite.} Let
$H = \mathsf{SHA256}(\pi)$. On its own $H$ is \emph{inert}. A prover can always
emit a self-consistent certificate for any policy $\pi'$ it prefers, simply by
placing $H' = \mathsf{SHA256}(\pi')$ in that certificate. \emph{Secrecy of
$\pi$ is what creates the gap}, and no cryptographic mechanism can close it:
closing it would require showing $\pi$ to the verifier, which is precisely what
the mode exists to avoid. What closes it is an \emph{issuing act}, so we state
it as an explicit prerequisite of the mode:

\begin{description}
  \item[(P-PP)] An \emph{issuing authority} signs $\pi$ and registers $H$ in a
  record readable by third parties. A policy is \emph{accountable} iff its
  \texttt{policy\_hash} can be traced to such a record.
\end{description}

Without (P-PP) the mode is not merely weaker---the verifier's conclusion is
empty. This is the same assumption as ``honest policy author'' in
\texttt{docs/security-model.md}~\S0.2; what is new here is scoping it to the
verification interface, which had not been done, and which a referee will
otherwise use to dismiss the whole mode in one sentence.

\paragraph{Three parties.} ``Three'' means \emph{authority / operator / third
party}; an auditor is a named instance of the third party, not a fourth trust
anchor.

\begin{table}[t]
\centering
\caption{The three-party model of policy-private mode.}
\label{tab:wp4-parties}
\setlength{\tabcolsep}{3pt}
\begin{tabular}{@{}L{0.15\columnwidth}L{0.13\columnwidth}L{0.64\columnwidth}@{}}
\toprule
Party & Honest? & Holds; wants \\
\midrule
Authority $R$ & \textbf{Yes}, by (P-PP) & Holds $\pi$ and the issuing record;
wants $\pi$ actually enforced, and wants to be able to check that it was. \\
Operator $O$ & \textbf{No} (the adversary) & Holds $\pi$, $T$, the proving key;
wants $\pi$ hidden from third parties while still being able to certify. \\
Third party $V$ & Yes & Holds only public values---\texttt{policy\_\allowbreak hash},
\texttt{circuit\_\allowbreak hash}, \texttt{response\_\allowbreak commitment}, \texttt{passed}, the
proof; wants to confirm that \emph{some issued} policy passed, and must
\emph{not} learn $\pi$. \\
\emph{(Auditor $V^{*}$)} & Yes & As $V$, plus $\pi$ once $O$ discloses it;
wants ex-post accountability. \\
\bottomrule
\end{tabular}
\end{table}

Two boundaries must be stated with the model. First, secrecy of $\pi$ comes
from a \emph{data-flow fact}---$\pi$ never enters the public values---not from
zero-knowledge; any rule that writes a fragment of $\pi$ into the public values
(the leaks in Table~\ref{tab:wp4}, items~3 and~5) breaks it on that fragment.
Second, accountability for $V^{*}$ is \emph{voluntary}: $V^{*}$ can check only
when $O$ chooses to disclose $\pi$. The mode offers no forced disclosure; that
would need the release-gating family discussed in
\texttt{docs/tdsc-plan.md}~\S1.6(b), which we do not claim.

\paragraph{Verifier capability table.} This table is the contribution; the
cryptography is downstream of it, and the table survives even if the
cryptography only partially lands. \emph{It must appear in the paper as
written}, because it is what stops \texttt{policy\_hash} from being read as
``compliance was proved''.

\begin{table}[t]
\centering
\caption{What each verifier may and may not conclude in policy-private mode.}
\label{tab:wp4-capability}
\footnotesize
\setlength{\tabcolsep}{3pt}
\begin{tabular}{@{}L{0.14\columnwidth}L{0.31\columnwidth}L{0.30\columnwidth}L{0.17\columnwidth}@{}}
\toprule
Verifier & \textbf{May} conclude & \textbf{May not} conclude & Rests on \\
\midrule
Third party $V$ & (i) the commitment binds $T$
(Definition~\ref{def:k}); (ii) \emph{some} policy with hash $H$ passed;
(iii) $H$ appears in the issuing record & (iv) the \emph{content} of $\pi$;
(v) how many rules $\pi$ has, or of which classes; (vi) ``the operator
enforces the regulation \emph{I} care about''---only ``some issued
regulation'' & (P-PP); A2; the $k$ of T1 \\
End user (also receiver) & As $V$, and additionally: the proved $T$ is the
string it received, since its nonce is bound (T3$'$, requires $f>0$) & Only
the string it received; \emph{not} what was delivered to another downstream
agent & $f>0$ \\
Authority $R$ & Everything; it holds $\pi$ and can recompute $H$ & --- & --- \\
Auditor $V^{*}$ (after disclosure) & (vii) $\mathsf{SHA256}(\pi) =
\texttt{policy\_\allowbreak hash}$, i.e.\ what ran was the issued $\pi$;
(viii) the disclosed violation record is consistent & (ix) the undisclosed
part of $T$ & voluntary disclosure \\
On-chain verifier & (i)--(iii) only; it hashes and verifies signatures, and
does not interpret policy semantics & (iv)--(ix) & (P-PP) \\
\bottomrule
\end{tabular}
\end{table}

The $V$ row is the whole reason this subsection exists: $V$ does not obtain
``compliance'', it obtains ``some issued policy was satisfied''. Conflating the
two is the sentence a referee would use to kill the mode.

\need{Remaining work on this subsection: (a) confirm the (P-PP) framing with
the collaborator, and choose its instantiation (an on-chain registry of issued
policy hashes vs.\ a signed manifest); (b) add the prior-art citations below;
(c) decide whether the capability table merges with Table~\ref{tab:wp4} in the
submission, or stays separate (currently separate, on the grounds that one
table is about leakage and the other about verifier conclusions).}

\subsection{Expected results, stated in advance}
\label{sec:wp4-expect}

\need{Calibrated expectations, so that a partial outcome is a publishable
result rather than a failure:}
\begin{itemize}
  \item Sub-problems 1, 2, 6: high confidence ($\gtrsim 80\%$). 1 is already
        true of the code; 2 follows from WP1; 6 is a design restriction.
  \item Sub-problems 3, 4, 5: likely ($\sim 70\%$).
  \item ``\emph{Zero}-knowledge violation evidence'' in the strong sense:
        only $50$--$60\%$. The boolean ``some violation occurred'' is easy; the
        hard part is \emph{granularity}---with $3$--$8$ rules, a coarsened
        violation class can itself identify the rule. \textbf{The likely best
        result is a granularity/actionability/leakage tradeoff curve, not full
        hiding}---and that curve is a contribution in its own right.
  \item \S\ref{sec:wp4-tension}, the tension and the capability table:
        $\sim 90\%$. This is precisely why the section leads with it.
\end{itemize}

\need{Prior-art discipline, and this is the highest-risk place in the paper ---
the collaborator's own field. \emph{Policy hiding is not a new primitive}:
hidden credentials, secret handshake, and XACML privacy in access control, and
attribute-hiding / fuzzy predicate encryption, are all the same family. Cite
that line explicitly and do not claim primacy. What is new is narrow: (a) the
\emph{application} to LLM output compliance together with selective disclosure
of violation type and location, and (b) the accountability tension of
\S\ref{sec:wp4-tension}.}

% =====================================================================
\section{Composition: The Constructive Half of the End-User-Secret Family}
\label{sec:composition}
% 决策 D4：双载体（ezkl + 真 ONNX 主臂；DeepProve 对照臂）。

\fixme{关键更正：组合本身不闭合缺口。只有当模型证明绑定到\emph{用户提交的
输入} $x$ 时，$M(x) = T \wedge \pi(T)$ 才真的闭合"送错字节"——而那恰好是
Theorem \ref{thm:t3} 要求的 $f > 0$：用户提交的 $x$ 就是那 $f$ 个新鲜比特。}

\need{组合契约的定义（模型证明须暴露什么，策略证明才能绑上）；
真模型规模（$\ge 10^6$ 参数）；成本数据。}

% =====================================================================
\section{Availability: The Unanchored-Pattern Cost Attack}
\label{sec:dos}
% WP6。

\begin{theorem}[Quadratic cost inflation]
\need{对含未锚定模式的正则，控制响应的对手可强制 $\Omega(n^2)$ 证明工作量。}
\end{theorem}

\need{缓解：Pike VM（已有）+ 响应长度承诺 + 分块匹配。
zk-regex 有 DFA 限制但没有 zkVM 场景下的 DoS 分析——这是新的。}

% =====================================================================
\section{Coverage: Can a Provider Simply Withhold Certificates?}
\label{sec:coverage}
% WP5。诚实预期：无可信硬件时可能无解 —— 那就写成 open problem。

% =====================================================================
\section{Evaluation}
\label{sec:eval}

\need{矩阵：规则 \{1,2,3,4,6,8\} $\times$ 长度 \{200, 2k, 10k, 50k\}
$\times$ 模式 \{transparent, private, private-policy\}
$\times$ 证明档 \{core, groth16\}。
骨架见 bench/bench_scale.py。}

\subsection{Why zero-knowledge? A decision map}
% WP9：非 ZK 基线 + 决策图。

\need{签名 Merkle 日志对照臂（brivora/VAOL 风格）。
结论：ZK 只在 \{私有响应, 私有策略\} 两个角上有正收益。
把"透明模式 ZK 冗余"从批评变成有刻度的设计结论。}

% =====================================================================
\section{Limitations}
\label{sec:limitations}

\need{必须写进去的：下列五条。每条都要在正文里展开，不是罗列。}
\begin{itemize}
  \item Content privacy here is \emph{commitment} privacy, not
        irrecoverability of content (Corollary \ref{cor:reexec}).
  \item The zero-knowledge property of the default proof mode:
        \fixme{SP1's published security model states explicitly that
        individual STARK proofs do not currently satisfy the
        zero-knowledge property. If groth16 does not become the delivery
        mode, the word ``zero-knowledge'' must leave the title.}
  \item Binding the composition's input requires the user to be online (or to
        hold a key). This is the price of Remark \ref{rem:taxonomy}(b), and
        the concrete form of $f > 0$ in Theorem \ref{thm:t3}.
  \item The coverage problem (\S\ref{sec:coverage}) may be unsolvable without
        trusted hardware.
  \item Three-party collusion is undetectable at the signature layer
        (\texttt{docs/security-model.md}).
\end{itemize}

% =====================================================================
\section{Conclusion}

% =====================================================================
\appendix
\section{Proofs}
\need{\textbf{状态（2026-10-12 自检后）}：
Theorem \ref{thm:t1}(H) 与 Theorem \ref{thm:t3} 的证明\textbf{已写在正文}，
此处不再重复。三分法已降级为分类 Remark \ref{rem:taxonomy}，无需证明。
Theorem \ref{thm:t1}(L) 无闭式，只给测量口径（\S\ref{sec:eval}）。
原 T2「分离定理」已\textbf{停用}（见 Remark \ref{rem:designated}）——
本节因此只剩一条待补事项，见下。}

\need{\textbf{待补}：把 Theorem \ref{thm:t1}(H) 的随机预言机模型换成
标准模型（或明确声明本文只在 ROM 下主张）。当前证明依赖 $H$ 是新鲜随机函数
这一条；实践中的 SHA-256 只给计算意义下的近似，故正文措辞必须写成
\emph{heuristic}，不得写成信息论界。}

% =====================================================================
\begin{thebibliography}{99}

\bibitem{hollowllm}
C.~Gong, B.~Liu, and M.~Li, ``Hollow-LLM Attack: Computationally Trivial
Weights in Zero-Knowledge Verification of LLM Inference,'' in \emph{IEEE S\&P},
2026, pp.~4339--4355, DOI 10.1109/SP63933.2026.00258.

\bibitem{zkattach}
\need{zk-attach, arXiv:2608.30083 --- 打开原文核实作者后再引}

\bibitem{zkgen}
J.~Lauinger, J.~Ernstberger, and S.~Steinhorst, ``zkGen: Policy-to-Circuit
Transpiler,'' in \emph{IEEE ICBC}, 2024, pp.~620--624,
DOI 10.1109/ICBC59979.2024.10634440.

\bibitem{pco}
\need{PCO, IACR ePrint 2026/994 --- 打开原文核实作者（Dolev/Yung 等）后再引}

\bibitem{zkagent}
L.~Wang, H.~Lou, C.~Li, Y.~Yu, and Y.~Hu, ``zkAgent: Verifiable LLM Agent
Execution via One-Shot Transcript Proofs,'' IACR ePrint 2026/199.

\bibitem{vdr}
\need{Verifiable Document Redacting, Springer, DOI 10.1007/978-3-319-66402-6_20
--- 核实作者与年份。\textbf{绝不能再引 IACR ePrint 2024/813}（那是
Kiraz/Larraia/Vaughan, "How to Redact the Bitcoin Backbone Protocol"）。
见 docs/tdsc-fixlist.md \S A。}

\end{thebibliography}

\end{document}
"   % 两次，第二遍收交叉引用
%  \SKELQUIET 让所有中文标记宏吞掉参数；不带它则需要能排版中文的引擎
%  （xelatex + CJK 字体），本骨架未配置。
% =====================================================================
\documentclass[journal]{IEEEtran}

\usepackage{amsmath,amssymb,amsthm}
\usepackage{booktabs}
\usepackage{array}
\usepackage{graphicx}
\usepackage{xcolor}
\usepackage{hyperref}

% 窄栏里用 p{} 会因两端对齐产生大片字距（编译时刷一堆 Underfull \hbox）。
% 表格正文列一律用 L{}：左对齐、允许自然参差。
\newcolumntype{L}[1]{>{\raggedright\arraybackslash}p{#1}}
\newcommand{\tabfont}{\footnotesize}

\newtheorem{theorem}{Theorem}
\newtheorem{lemma}[theorem]{Lemma}
\newtheorem{proposition}[theorem]{Proposition}
\newtheorem{corollary}[theorem]{Corollary}
\newtheorem{definition}[theorem]{Definition}
\newtheorem{assumption}[theorem]{Assumption}
\newtheorem{remark}[theorem]{Remark}

% --- 工作标记：提交前全部清掉 ---
% 注意：宏名不能叫 \wp —— \wp 已是 LaTeX/amssymb 的 Weierstrass ℘，
% \newcommand{\wp} 会直接报 "Command \wp already defined" 并使后文级联报错。
\ifdefined\SKELQUIET
  \newcommand{\wpmark}[1]{}
  \newcommand{\need}[1]{}
  \newcommand{\fixme}[1]{}
\else
  \newcommand{\wpmark}[1]{\textcolor{red}{\textsuperscript{[WP#1]}}}
  \newcommand{\need}[1]{\textcolor{blue}{[[ NEED: #1 ]]}}
  \newcommand{\fixme}[1]{\textcolor{orange}{[[ FIX: #1 ]]}}
\fi

\begin{document}

% ---------------------------------------------------------------------
% 标题：决策 D5 —— 保留 "Zero-Knowledge" 的前提是 groth16 成为交付档 (WP1)。
%       若 WP1 未打通 groth16，必须删掉标题与摘要里的 zero-knowledge。
% ---------------------------------------------------------------------
\title{Proof-of-Policy: Zero-Knowledge Compliance Certificates\\ for LLM Agent Outputs}
% 备选标题（若 groth16 未打通）："Proof-of-Policy: Succinct Compliance
% Certificates for LLM Agent Outputs"

\author{\need{authors}}
\maketitle

% ---------------------------------------------------------------------
\begin{abstract}
% 决策 E2：内容隐私上界必须出现在摘要，不能埋进 §8.1。
We present \emph{Proof-of-Policy} (PoP), a system that proves an LLM agent's
response $T$ satisfies a publicly declared policy $\pi$, without a trusted
execution environment...
\need{一句完整的贡献陈述（对应 D1 的"分析 x 系统各半"）}

\smallskip
\noindent\textbf{Scope of privacy.} PoP provides \emph{policy}-zero-knowledge:
the certificate reveals neither the policy evaluation trace nor, in
policy-private mode, the policy itself. It does \emph{not} guarantee that $T$
is unrecoverable: as we show in \S\ref{sec:binding}, any publicly recomputable
binding to $T$ admits a re-execution attack whenever the adversary knows the
prompt and the model. The same bit-budget governs the request side: a
certificate is replayable without limit unless the receiver contributes fresh
randomness (Theorem \ref{thm:t3}).
\end{abstract}

\begin{IEEEkeywords}
zero-knowledge proofs, LLM agents, policy compliance, verifiable computation,
accountability
\end{IEEEkeywords}

% =====================================================================
\section{Introduction}
\label{sec:intro}

\need{Contribution list. 必须包含下面这条"不主张"声明：}

\noindent\textbf{What we do \emph{not} claim.}
We do not claim that policy proofs bind nothing: our own system already binds
the response via a verifier-issued nonce (Definition 1). We do not claim to be
the first to observe that a proof can fail to bind the artifact it purports to
describe: the \emph{effort gap} for ZK inference proofs is formalised by
Hollow-LLM \cite{hollowllm}, and the \emph{source-integrity gap} for predicate
proofs is named by zk-attach \cite{zkattach}. Our contribution is to
(i) show that the surviving gap is the \emph{divergent-delivery} one
(Proposition \ref{thm:gap}), (ii) give a quantitative bound on what a public
value reveals, together with the \emph{structural hypothesis} under which that
bound holds---and a demonstration that the hypothesis is necessary, since our
own two public values sit on opposite sides of it (Theorem \ref{thm:t1} and
Remark \ref{rem:betanecessary}), and (iii) give a quantitative price for
closing the gap (Theorem \ref{thm:t3}).

\smallskip
\noindent\textbf{A separation we deliberately do \emph{not} claim.} It is
tempting to promote our designated-verifier interface into a theorem
(``public verification is impossible when $k$ is small''). We give the reason
we do not in Remark \ref{rem:designated}: in the low-$k$ regime hiding fails
for every scheme, designated or not, and the dichotomy has no content of its
own. We state this explicitly because the negated claim is the more attractive
one to write.

\smallskip
\noindent\textbf{One spine, two sides.} Theorems \ref{thm:t1} and
\ref{thm:t3} state the same fact about two different inputs: \emph{a proof
binds only as much as the committed, adversary-non-selectable quantity has bits
for}. The response side is governed by entropy,
$k = -\log \mathsf{Guess}(T \mid Z)$ (Definition \ref{def:k}); the request side
by freshness, $f$. Neither can be bought with more computation.

\fixme{全文禁止出现"比 Proof-of-Guardrail 更强"或任何强弱对比。
Proof-of-Guardrail 与本文正交：它证"计算发生过"，本文证"性质成立"。}

% =====================================================================
\section{Background and Related Work}
\label{sec:related}

\need{五条线，每条都要"它占什么位 / 我们差在哪"两列。见 docs/tdsc-plan.md \S2。}

\begin{table}[t]
\centering
\caption{Positioning against the five prior-art lineages.}
\label{tab:related}
\setlength{\tabcolsep}{3pt}
\begin{tabular}{@{}L{0.26\columnwidth}L{0.42\columnwidth}L{0.26\columnwidth}@{}}
\toprule
Lineage & Representative work & Occupied slot \\
\midrule
Policy $\to$ circuit      & zkGen \cite{zkgen}, Circuitree, Prezta & DSL-to-circuit \\
Agent/LLM compliance      & NiyamAI, Zaguk, PCO \cite{pco}, WITNESS & private-mode kernel \\
Model computation         & DeepProve, zkGPT, zkAgent \cite{zkagent} & inference integrity \\
Attacks / negative results& Hollow-LLM \cite{hollowllm}, zk-attach \cite{zkattach} & naming the gap \\
Non-ZK industry baseline  & @brivora/verify, VAOL & signed audit logs \\
\bottomrule
\end{tabular}
\end{table}

\fixme{引文纪律：除 Hollow-LLM / zkGen / NiyamAI / zkAgent / ePrint 2024/813
五条外，其余引文必须逐条打开原文核实后再进 bibliography。
见 docs/tdsc-fixlist.md \S2 的注记。}

% =====================================================================
\section{System Overview}
\label{sec:system}
% 决策 D1 的"系统半边"：这一节压缩到 ~1.5 页，是 T3(b) 的载体，不是主角。

\need{DSL $\to$ ConstraintSpec $\to$ circuit；双模式；证书 + 锚定；
工具回执链；801 测试 / 15 skip 的工件陈述。}

% =====================================================================
\section{Binding of Policy Proofs}
\label{sec:binding}
% 决策 D1 的"分析半边"：这是论文的核心。

\subsection{The gap}

\begin{definition}[Policy-proof system]
A policy-proof system is a tuple $(\mathsf{Prove}, \mathsf{Ver}, \pi, \mathcal{M})$.
On challenge $n$ from a verifier, $\mathsf{Prove}(\pi, T, n)$ outputs a
certificate $c$ whose public values include
\begin{multline*}
b \;=\; H(\texttt{pop-bind-v1} \,\|\, \mathsf{u32be}(|n|) \,\|\, n \,\|\, T),\\
C = H(T).
\end{multline*}
$\mathsf{Ver}(c, n, T')$ accepts iff the policy relation holds \emph{and}
$b$ recomputed from the received $T'$ matches.
\end{definition}

\fixme{\textbf{H3/H4}：早先写的 ``the relation contains no term referring to the
delivered bytes'' 与仓库事实冲突——公开值里\emph{已有} $b$，验证时由送达的
$T'$ 现场重算（\texttt{circuits/types/src/lib.rs:444}、
\texttt{docs/security-model.md:232}）。命题已按上面的定义改写。见
\texttt{docs/tdsc-fixlist.md} \S H。}

\begin{proposition}[Divergent delivery]
\label{thm:gap}
Definition 1 binds $T'$ \emph{as received by the verifying party}. If the
provider delivers a different string $T'' \ne T'$ to the end user, verification
still succeeds on $T'$. Neither hardware nor a proof can prevent this; only the
end user holding $n$ and recomputing $b$ can.
\end{proposition}
\begin{proof}
$b$ is a function of $(n,T)$ only; the delivery channel to the end user carries
no authenticated term. One line.
\end{proof}

\need{This is the \emph{residual} gap, not ``proofs bind nothing''. Prior work
must be cited here (Hollow-LLM \cite{hollowllm}, zk-attach \cite{zkattach}), and
we state plainly that the naive claim would be false against our own system.}

\subsection{A quantitative binding--hiding tradeoff}
\label{sec:tradeoff}

% ---------------------------------------------------------------------
% T1 : 关键是把"对手已知信息 Z"定对。
%      docs/tdsc-plan.md §1.4 已修正原 §4 的非因果推理。
% ---------------------------------------------------------------------
\begin{definition}[Guessing entropy of a response]
\label{def:k}
For a response distribution $\mathcal{D}$ and adversary side information $Z$,
let
\[
\mathsf{Guess}(T \mid Z) \;\triangleq\; \max_{\mathcal{A}} \Pr[\mathcal{A}(Z) = T],
\]
which equals $\mathbb{E}_z[\max_t \Pr[T = t \mid Z = z]]$, and
\[
k \;\triangleq\; -\log_2 \mathsf{Guess}(T \mid Z),
\]
the maximum ranging over all time-unbounded adversaries.
\end{definition}

\begin{remark}[The definition is operational, and it is the average]
\label{rem:kmeas}
$k$ is \emph{measured}, not assumed: sampling one model on one prompt many
times estimates $\mathsf{Guess}(T \mid Z)$ directly
(\S\ref{sec:eval}), so Theorem \ref{thm:t1} is instantiated with a
number we obtain rather than a number we posit. The definition must average
over $z$ rather than take the pointwise worst case
$k_z = -\log \max_t \Pr[T = t \mid Z = z]$: the proof takes an expectation over
$z$, and $\mathbb{E}_z[2^{-k_z}] = 2^{-k}$ holds only for the average form.
\end{remark}

\begin{assumption}[Adversary side information]
\label{as:Z}
The adversary knows the prompt, the model (possibly its weights), and the
decoding parameters. \emph{This is the assumption that makes the theorem
non-trivial}: under greedy decoding, $k \approx 0$.
\end{assumption}

\begin{theorem}[Binding, and the hypothesis it needs]
\label{thm:t1}
Let $\beta : \mathcal{M} \to \{0,1\}^\lambda$ be public, deterministic and
efficiently computable, and let the certificate publish $\beta(T)$. Let
$\mathcal{A}$ be any adversary, unbounded in time, that is given $Z$ and
evaluates $\beta$ at most $q$ times.
\begin{enumerate}
\item[(H)] \emph{(Randomised public value.)} Suppose
  $\beta(T) = H(\mathsf{dom} \,\|\, \mathsf{u32be}(|n|) \,\|\, n \,\|\, T)$ for a
  random oracle $H$ and a nonce $n$ drawn uniformly and \emph{independently of
  $T$}. Then
  \[
  \Pr\!\left[\mathcal{A}^\beta\!\left(Z, \beta(T)\right) = T\right]
  \;\le\; q \cdot 2^{-k} \;+\; q \cdot 2^{-\lambda}.
  \]
  Consequently hiding with advantage $\varepsilon$ requires
  $k \ge \log q + \log(1/\varepsilon)$.
\item[(L)] \emph{(Lossy public value.)} For an arbitrary deterministic public
  $\beta$ the correct bound is over the \emph{preimage class}:
  \begin{multline*}
  \Pr\!\left[\mathcal{A}^\beta\!\left(Z, \beta(T)\right) = T\right]
  \;\le\; \mathsf{Guess}\!\left(T \mid Z, \beta(T)\right) \\
  \triangleq\; \mathbb{E}_y\!\left[\max_t \Pr[T = t \mid Z, \beta(T) = y]\right].
  \end{multline*}
\end{enumerate}
\end{theorem}

\begin{proof}
\emph{(H).} Fix $z$ and let the adversary's query set be
$Q = \{q_1, \dots, q_q\}$. Decompose by first hit,
\[
\Pr[\exists i : q_i = T] \;=\; \sum_i \Pr\!\left[q_i = T \wedge \neg E_{<i}\right],
\]
where $E_{<i} \triangleq \{\exists j < i : q_j = T\}$.
Conditioned on $\neg E_{<i}$, every oracle answer $a_1,\dots,a_{i-1}$ already
revealed is uniform and independent of $T$ (no queried point equals $T$, and $H$
is a fresh random function), and $\beta(T)$ itself is uniform and independent of
$T$. Hence $q_i$---a deterministic function of $(z, \beta(T), a_{<i})$---is
\emph{independent of $T$}, so
\begin{multline*}
\Pr[q_i = T \mid Z = z]
= \sum_t \Pr[q_i = t]\,\Pr[T = t \mid Z = z] \\
\le \max_t \Pr[T = t \mid Z = z] = 2^{-k_z},
\end{multline*}
giving $\Pr[\exists i : q_i = T \mid z] \le q\,2^{-k_z}$. Taking the expectation
over $z$ and using $\mathbb{E}_z[2^{-k_z}] = \mathsf{Guess}(T \mid Z) = 2^{-k}$
from Definition \ref{def:k} yields $q \cdot 2^{-k}$. An adversary that outputs
an unqueried value wins with probability at most $2^{-k_z} \le q\,2^{-k_z}$,
which the same bound absorbs. The term $q \cdot 2^{-\lambda}$ charges the
probability that a query hits the uniform value $\beta(T)$ by chance.

\emph{(L).} Immediate from Definition \ref{def:k} after conditioning on
$\beta(T)$; no hypothesis on $\beta$ is used.
\end{proof}

\begin{remark}[The hypothesis in (H) is necessary, and our two instances straddle it]
\label{rem:betanecessary}
Without the hypothesis, part (H) is not merely unproven but \emph{false}: take
$\beta = \mathrm{id}$, or $\beta = \mathsf{encode}$ as in
Proposition \ref{prop:semantic}. The adversary outputs $\beta(T)$ and succeeds
with probability $1$, independently of $q$ and $k$. Our own two public values
therefore sit on opposite sides of the hypothesis---$b =
H(\mathsf{dom}\|n\|T)$ satisfies it, $\mathsf{encode}(T)$ does not---and this is
what makes the pair a theorem about \emph{our design} rather than a citation of
a standard bound. An earlier draft of this work stated (H) without its
hypothesis; the counterexample above is why the hypothesis is now explicit.
\end{remark}

\begin{corollary}[Re-execution, not enumeration]
\label{cor:reexec}
Under Assumption \ref{as:Z}, when the decoding is greedy or low-temperature,
$k \approx 0$ and the attack of Theorem \ref{thm:t1} is realised by
\emph{re-running the model}, at the cost of one inference---independent of
$|T|$ and requiring \emph{no} evaluation of $\beta$. Enumeration over the string
space is \emph{not} the relevant attack.
\end{corollary}

\fixme{原 docs/sp1-zk-audit.md §4 写的是"熵(~10^2--10^3 bit)远低于 256 bit，
因此枚举可行"——这是非因果推论（2^100 与 2^256 一样枚举不了；真正的攻击是\textbf{重执行}，
对手的已知信息是 (prompt, 模型, 解码参数)）。
2026-10-12 已全库改正，docs/tdsc-fixlist.md §G 的 G1--G6 逐条结案。}

\begin{remark}[A design observation, not a separation]
\label{rem:designated}
Our verification interface is \emph{designated}: binding is checked against a
value $b = H(\mathsf{dom}\|n\|T)$ in which a verifier-issued nonce $n$ appears.
It is tempting to promote this into a separation (``public verification is
impossible when $k$ is small''). \textbf{We claim no such theorem, and
Corollary \ref{cor:reexec} is why.} When $k \approx 0$, hiding fails for
\emph{every} scheme, designated or not, because the adversary re-runs the
model; when $k$ is large, a zero-knowledge proof of ``$\exists\, T : b =
H(\mathsf{dom}\|n\|T) \wedge J(\pi, T) = 1$'' is simultaneously publicly
verifiable and hiding. The dichotomy is thus governed by $k$ alone and carries
no independent content. What does survive---and what
\S\ref{sec:wp4-tension}
develops---is that \emph{third-party auditing of delivery} forces disclosure of
$T$, which is an informational remark rather than a theorem. We record the
demotion here because an earlier draft asserted the separation as a theorem.
\end{remark}

\need{\textbf{This is retrospective justification, not a novel construction.}
The system publishes $b = H(\mathsf{dom}\|n\|T)$ with a verifier-issued nonce
$n$, and Remark \ref{rem:designated} explains the shape of that choice. The
security model does not pin down whether $n$ is public; making it public
forfeits hiding, and making it verifier-held is what makes verification
designated.}

\need{\textbf{实验内容}：实测 $k$。对一组真实 prompt、同一模型多次采样，
估 $\mathsf{Guess}(T\mid Z)$（Definition \ref{def:k}），代入
Theorem \ref{thm:t1}(H) 得阈值。
据先例扫查，\emph{没有人测过这个量}——它本身是一块贡献。}

\begin{proposition}[The lossy instance of Theorem \ref{thm:t1}(L)]
\label{prop:semantic}
For the semantic rule class, the public values contain $\beta(T) =
\mathsf{encode}(T)$, a content-preserving feature vector of the response. This
$\beta$ is deterministic and injective on the features but not on strings, so
it violates the hypothesis of Theorem \ref{thm:t1}(H) and instantiates part
(L): the binding force is exactly
$\mathsf{Guess}(T \mid Z, \mathsf{encode}(T))$, i.e.\ the guessing probability
\emph{within the preimage class} of the published vector. Here $k$ need not be
argued from perplexity at all: the adversary's search space is that preimage,
which is dramatically smaller than the string space. The relevant asymmetry is
not ``low entropy'' but a \emph{lossy public encoding}.
\end{proposition}
\begin{proof}
Immediate. The measurement is direct: publish $\mathsf{encode}(T)$ for real
responses and measure how much of $T$ is recoverable.
\end{proof}

\need{This proposition is the paper's cheapest experiment and its strongest
anchor for Theorem \ref{thm:t1}---it is measurable in a day, and it makes the
bound concrete rather than cited. Source: \texttt{docs/design-semantic-rules.md}
lines 195--201, where the repository itself already names it as a second
instance.}

\subsection{Quantitative request binding}
\label{sec:t3}
% 决策 D2（第二版，2026-10-12）：原「闭合三分法」定理降级为 Remark（分类），
% 取代它的是下面的定量界 T3′。理由见 docs/tdsc-plan.md §1.6(b)：
% 三分法要么假（理想功能模型下能构造出第四条路），要么琐碎（把三条路写成定义）。

\begin{theorem}[Selective-presentation advantage]
\label{thm:t3}
Fix a policy $\pi$ and a request distribution. A selective-presentation
adversary $\mathcal{B}$ may, \emph{before} seeing the receiver's challenge,
compute a set $S$ of accepted transcripts. The receiver then contributes $f$
fresh, uniformly random bits $r$ that enter the proven statement. Then
\[
\Pr\!\left[\mathcal{B}\ \text{is accepted on } r\right] \;\le\; |S| \cdot 2^{-f}.
\]
If instead the receiver contributes no fresh randomness ($f = 0$; e.g.\ an
empty nonce), the advantage is $1$, and \emph{does not decay in any parameter
of the system}.
\end{theorem}
\begin{proof}
\need{主论证：$r$ 进入 $b = H(\text{dom} \| r \| T)$，故 $S$ 里任何一个
预先算好的证书，只有在 $r$ 恰好落进它已覆盖的挑战集合时才被接受；
$f$ 个新鲜均匀比特给出 $2^{-f}$ 的对撞上界，再对 $S$ 取并界。
$f = 0$ 时 $S$ 中每个元素都是永久有效的重放。}

\need{\textbf{代码依据}：\texttt{circuits/types/src/lib.rs}:950--995 已写明
“nonce 为空是合法的……不提供任何重放防护”。定理的 $f = 0$ 分支不是假想，
是本系统当前 API 的\emph{合法输入}。}
\end{proof}

\need{这条定理与 Theorem \ref{thm:t1} \textbf{同构}，而正是这一点使它值得写：
两条说的是同一件事——\emph{一个证明只能绑住提交进去的、对手无法预先选择的
那部分比特所能支撑的量}。Theorem \ref{thm:t1} 管响应侧
（$k = -\log_2 \mathsf{Guess}(T \mid Z)$，见 Definition \ref{def:k}），Theorem \ref{thm:t3} 管请求侧（新鲜度 $f$）。
这条统一主线写在 \S\ref{sec:intro} 的贡献列表里，是全文的理论骨架。}

\begin{remark}[A taxonomy, not a trichotomy]
\label{rem:taxonomy}
Every policy-proof deployment we surveyed closes the divergent-delivery gap of
Proposition \ref{thm:gap} in one of three ways:
(a) root trust in hardware (attested enclaves);
(b) have the \emph{end user} contribute a secret to the proven statement (the
user issues the nonce); or
(c) gate delivery on successful verification (\emph{prove-before-deliver}).
\textbf{We offer this as a classification, not as an exhaustiveness theorem}:
we neither prove that a fourth family cannot exist, nor that the three are
disjoint. zk-attach \cite{zkattach} walks (a); NiyamAI and ZKAP walk (c); PoP
walks (b).
\end{remark}

\need{\textbf{Posture, stated plainly.} Our system already carries
\texttt{response\_binding}, but $n$ is issued by the \emph{verifier}. Moving it
to the \emph{end user}---making $f > 0$ the receiver's contribution---closes
the divergent-delivery branch of Proposition \ref{thm:gap}. This is a small
protocol amendment, not a new system, and the paper must say exactly that. What
Theorem \ref{thm:t3} buys is the \emph{price}: with $f = 0$ there is no closure
at any cost.}

% =====================================================================
\section{Policy-Private Compliance}
\label{sec:private-policy}
% 决策 D3：这是"透明模式 ZK 冗余"这个批评的唯一硬解。
% 2026-10-12 读代码后的定性：这不是"造一个新证明系统"，是一次\emph{披露策略}
% 的重构——策略文本在架构上已经是证人。见 docs/tdsc-plan.md WP4.1。

\need{动机：对确定性规则 + 公开响应，签名哈希链就够了。唯一无法用签名
替代的场景是策略本身要保密。}

\subsection{The observation that reframes this work package}
\label{sec:wp4-finding}

\texttt{spec\_canonical} is already a \emph{guest input}---a witness---not a
public value: only
$\mathsf{policy\_hash} = \mathsf{SHA256}(\texttt{spec\_canonical})$
enters the public values (\texttt{circuits/types/src/lib.rs}:1--16, 415--416,
918, 1003--1004). So ``do not reveal the policy text'' is
\emph{architecturally already true}; what leaks is everything \emph{derived}
from the policy. WP4 is accordingly a \emph{disclosure-policy refactor}, not a
new proof system, and must be written that way---as a systematic account of
what a certificate leaks about the policy it certifies.

\begin{table}[t]
\centering
\caption{What a policy-private certificate leaks today, and what closing each
leak costs.}
\label{tab:wp4}
\begin{tabular}{@{}L{0.04\linewidth}L{0.46\linewidth}L{0.36\linewidth}@{}}
\toprule
\# & Leak & Difficulty / dependency \\
\midrule
1 & The policy text itself & \textbf{Already closed} (it is a witness). \\
2 & Witness hiding in the proof system & Medium; \emph{inherited from WP1}:
    only \texttt{groth16}/\texttt{plonk} is witness-hiding, per
    \texttt{docs/\allowbreak sp1-zk-audit.md}. \\
3 & \texttt{PrivateViolation}\allowbreak\{\texttt{rule},\allowbreak\texttt{kind}\} names the rule that
    fired & Medium; needs a coarsened violation class plus a commitment to it. \\
4 & Evidence opening reveals the rule
    (\texttt{policydsl/\allowbreak privacy/\allowbreak commit.py}:319) & Medium; disclose the
    \emph{existence} of a violation, not its identity. \\
5 & The \texttt{delegated} flag exposes policy structure (which rules were
    routed to the ezkl path) & Medium; easiest to overlook, cheapest to fix. \\
6 & Semantic rules publish $\beta(T)=\mathsf{encode}(T)$
    (Proposition \ref{prop:semantic}) & Medium--high; either bound the
    preimage, or disable this rule class in policy-private mode. \\
\bottomrule
\end{tabular}
\end{table}

\need{Honest structure of the table, and it should be stated as such:
sub-problem 1 is a one-paragraph re-reading of the existing code, \emph{not} a
contribution; 2 is inherited wholesale from WP1; 3--5 are engineering with one
real formal core (disclosure granularity); 6 is where the work is---and it
collides with Theorem \ref{thm:t1}, because one cannot hide the policy while
publishing a lossy encoding of the response that the semantic rules require.}

\subsection{The tension worth writing up: policy privacy vs.\ accountability}
\label{sec:wp4-tension}

Lead the section with the tension, not with the cryptography: it determines
whether the mode means anything at all.

\paragraph{The prerequisite, stated as a prerequisite.} Let
$H = \mathsf{SHA256}(\pi)$. On its own $H$ is \emph{inert}. A prover can always
emit a self-consistent certificate for any policy $\pi'$ it prefers, simply by
placing $H' = \mathsf{SHA256}(\pi')$ in that certificate. \emph{Secrecy of
$\pi$ is what creates the gap}, and no cryptographic mechanism can close it:
closing it would require showing $\pi$ to the verifier, which is precisely what
the mode exists to avoid. What closes it is an \emph{issuing act}, so we state
it as an explicit prerequisite of the mode:

\begin{description}
  \item[(P-PP)] An \emph{issuing authority} signs $\pi$ and registers $H$ in a
  record readable by third parties. A policy is \emph{accountable} iff its
  \texttt{policy\_hash} can be traced to such a record.
\end{description}

Without (P-PP) the mode is not merely weaker---the verifier's conclusion is
empty. This is the same assumption as ``honest policy author'' in
\texttt{docs/security-model.md}~\S0.2; what is new here is scoping it to the
verification interface, which had not been done, and which a referee will
otherwise use to dismiss the whole mode in one sentence.

\paragraph{Three parties.} ``Three'' means \emph{authority / operator / third
party}; an auditor is a named instance of the third party, not a fourth trust
anchor.

\begin{table}[t]
\centering
\caption{The three-party model of policy-private mode.}
\label{tab:wp4-parties}
\setlength{\tabcolsep}{3pt}
\begin{tabular}{@{}L{0.15\columnwidth}L{0.13\columnwidth}L{0.64\columnwidth}@{}}
\toprule
Party & Honest? & Holds; wants \\
\midrule
Authority $R$ & \textbf{Yes}, by (P-PP) & Holds $\pi$ and the issuing record;
wants $\pi$ actually enforced, and wants to be able to check that it was. \\
Operator $O$ & \textbf{No} (the adversary) & Holds $\pi$, $T$, the proving key;
wants $\pi$ hidden from third parties while still being able to certify. \\
Third party $V$ & Yes & Holds only public values---\texttt{policy\_\allowbreak hash},
\texttt{circuit\_\allowbreak hash}, \texttt{response\_\allowbreak commitment}, \texttt{passed}, the
proof; wants to confirm that \emph{some issued} policy passed, and must
\emph{not} learn $\pi$. \\
\emph{(Auditor $V^{*}$)} & Yes & As $V$, plus $\pi$ once $O$ discloses it;
wants ex-post accountability. \\
\bottomrule
\end{tabular}
\end{table}

Two boundaries must be stated with the model. First, secrecy of $\pi$ comes
from a \emph{data-flow fact}---$\pi$ never enters the public values---not from
zero-knowledge; any rule that writes a fragment of $\pi$ into the public values
(the leaks in Table~\ref{tab:wp4}, items~3 and~5) breaks it on that fragment.
Second, accountability for $V^{*}$ is \emph{voluntary}: $V^{*}$ can check only
when $O$ chooses to disclose $\pi$. The mode offers no forced disclosure; that
would need the release-gating family discussed in
\texttt{docs/tdsc-plan.md}~\S1.6(b), which we do not claim.

\paragraph{Verifier capability table.} This table is the contribution; the
cryptography is downstream of it, and the table survives even if the
cryptography only partially lands. \emph{It must appear in the paper as
written}, because it is what stops \texttt{policy\_hash} from being read as
``compliance was proved''.

\begin{table}[t]
\centering
\caption{What each verifier may and may not conclude in policy-private mode.}
\label{tab:wp4-capability}
\footnotesize
\setlength{\tabcolsep}{3pt}
\begin{tabular}{@{}L{0.14\columnwidth}L{0.31\columnwidth}L{0.30\columnwidth}L{0.17\columnwidth}@{}}
\toprule
Verifier & \textbf{May} conclude & \textbf{May not} conclude & Rests on \\
\midrule
Third party $V$ & (i) the commitment binds $T$
(Definition~\ref{def:k}); (ii) \emph{some} policy with hash $H$ passed;
(iii) $H$ appears in the issuing record & (iv) the \emph{content} of $\pi$;
(v) how many rules $\pi$ has, or of which classes; (vi) ``the operator
enforces the regulation \emph{I} care about''---only ``some issued
regulation'' & (P-PP); A2; the $k$ of T1 \\
End user (also receiver) & As $V$, and additionally: the proved $T$ is the
string it received, since its nonce is bound (T3$'$, requires $f>0$) & Only
the string it received; \emph{not} what was delivered to another downstream
agent & $f>0$ \\
Authority $R$ & Everything; it holds $\pi$ and can recompute $H$ & --- & --- \\
Auditor $V^{*}$ (after disclosure) & (vii) $\mathsf{SHA256}(\pi) =
\texttt{policy\_\allowbreak hash}$, i.e.\ what ran was the issued $\pi$;
(viii) the disclosed violation record is consistent & (ix) the undisclosed
part of $T$ & voluntary disclosure \\
On-chain verifier & (i)--(iii) only; it hashes and verifies signatures, and
does not interpret policy semantics & (iv)--(ix) & (P-PP) \\
\bottomrule
\end{tabular}
\end{table}

The $V$ row is the whole reason this subsection exists: $V$ does not obtain
``compliance'', it obtains ``some issued policy was satisfied''. Conflating the
two is the sentence a referee would use to kill the mode.

\need{Remaining work on this subsection: (a) confirm the (P-PP) framing with
the collaborator, and choose its instantiation (an on-chain registry of issued
policy hashes vs.\ a signed manifest); (b) add the prior-art citations below;
(c) decide whether the capability table merges with Table~\ref{tab:wp4} in the
submission, or stays separate (currently separate, on the grounds that one
table is about leakage and the other about verifier conclusions).}

\subsection{Expected results, stated in advance}
\label{sec:wp4-expect}

\need{Calibrated expectations, so that a partial outcome is a publishable
result rather than a failure:}
\begin{itemize}
  \item Sub-problems 1, 2, 6: high confidence ($\gtrsim 80\%$). 1 is already
        true of the code; 2 follows from WP1; 6 is a design restriction.
  \item Sub-problems 3, 4, 5: likely ($\sim 70\%$).
  \item ``\emph{Zero}-knowledge violation evidence'' in the strong sense:
        only $50$--$60\%$. The boolean ``some violation occurred'' is easy; the
        hard part is \emph{granularity}---with $3$--$8$ rules, a coarsened
        violation class can itself identify the rule. \textbf{The likely best
        result is a granularity/actionability/leakage tradeoff curve, not full
        hiding}---and that curve is a contribution in its own right.
  \item \S\ref{sec:wp4-tension}, the tension and the capability table:
        $\sim 90\%$. This is precisely why the section leads with it.
\end{itemize}

\need{Prior-art discipline, and this is the highest-risk place in the paper ---
the collaborator's own field. \emph{Policy hiding is not a new primitive}:
hidden credentials, secret handshake, and XACML privacy in access control, and
attribute-hiding / fuzzy predicate encryption, are all the same family. Cite
that line explicitly and do not claim primacy. What is new is narrow: (a) the
\emph{application} to LLM output compliance together with selective disclosure
of violation type and location, and (b) the accountability tension of
\S\ref{sec:wp4-tension}.}

% =====================================================================
\section{Composition: The Constructive Half of the End-User-Secret Family}
\label{sec:composition}
% 决策 D4：双载体（ezkl + 真 ONNX 主臂；DeepProve 对照臂）。

\fixme{关键更正：组合本身不闭合缺口。只有当模型证明绑定到\emph{用户提交的
输入} $x$ 时，$M(x) = T \wedge \pi(T)$ 才真的闭合"送错字节"——而那恰好是
Theorem \ref{thm:t3} 要求的 $f > 0$：用户提交的 $x$ 就是那 $f$ 个新鲜比特。}

\need{组合契约的定义（模型证明须暴露什么，策略证明才能绑上）；
真模型规模（$\ge 10^6$ 参数）；成本数据。}

% =====================================================================
\section{Availability: The Unanchored-Pattern Cost Attack}
\label{sec:dos}
% WP6。

\begin{theorem}[Quadratic cost inflation]
\need{对含未锚定模式的正则，控制响应的对手可强制 $\Omega(n^2)$ 证明工作量。}
\end{theorem}

\need{缓解：Pike VM（已有）+ 响应长度承诺 + 分块匹配。
zk-regex 有 DFA 限制但没有 zkVM 场景下的 DoS 分析——这是新的。}

% =====================================================================
\section{Coverage: Can a Provider Simply Withhold Certificates?}
\label{sec:coverage}
% WP5。诚实预期：无可信硬件时可能无解 —— 那就写成 open problem。

% =====================================================================
\section{Evaluation}
\label{sec:eval}

\need{矩阵：规则 \{1,2,3,4,6,8\} $\times$ 长度 \{200, 2k, 10k, 50k\}
$\times$ 模式 \{transparent, private, private-policy\}
$\times$ 证明档 \{core, groth16\}。
骨架见 bench/bench_scale.py。}

\subsection{Why zero-knowledge? A decision map}
% WP9：非 ZK 基线 + 决策图。

\need{签名 Merkle 日志对照臂（brivora/VAOL 风格）。
结论：ZK 只在 \{私有响应, 私有策略\} 两个角上有正收益。
把"透明模式 ZK 冗余"从批评变成有刻度的设计结论。}

% =====================================================================
\section{Limitations}
\label{sec:limitations}

\need{必须写进去的：下列五条。每条都要在正文里展开，不是罗列。}
\begin{itemize}
  \item Content privacy here is \emph{commitment} privacy, not
        irrecoverability of content (Corollary \ref{cor:reexec}).
  \item The zero-knowledge property of the default proof mode:
        \fixme{SP1's published security model states explicitly that
        individual STARK proofs do not currently satisfy the
        zero-knowledge property. If groth16 does not become the delivery
        mode, the word ``zero-knowledge'' must leave the title.}
  \item Binding the composition's input requires the user to be online (or to
        hold a key). This is the price of Remark \ref{rem:taxonomy}(b), and
        the concrete form of $f > 0$ in Theorem \ref{thm:t3}.
  \item The coverage problem (\S\ref{sec:coverage}) may be unsolvable without
        trusted hardware.
  \item Three-party collusion is undetectable at the signature layer
        (\texttt{docs/security-model.md}).
\end{itemize}

% =====================================================================
\section{Conclusion}

% =====================================================================
\appendix
\section{Proofs}
\need{\textbf{状态（2026-10-12 自检后）}：
Theorem \ref{thm:t1}(H) 与 Theorem \ref{thm:t3} 的证明\textbf{已写在正文}，
此处不再重复。三分法已降级为分类 Remark \ref{rem:taxonomy}，无需证明。
Theorem \ref{thm:t1}(L) 无闭式，只给测量口径（\S\ref{sec:eval}）。
原 T2「分离定理」已\textbf{停用}（见 Remark \ref{rem:designated}）——
本节因此只剩一条待补事项，见下。}

\need{\textbf{待补}：把 Theorem \ref{thm:t1}(H) 的随机预言机模型换成
标准模型（或明确声明本文只在 ROM 下主张）。当前证明依赖 $H$ 是新鲜随机函数
这一条；实践中的 SHA-256 只给计算意义下的近似，故正文措辞必须写成
\emph{heuristic}，不得写成信息论界。}

% =====================================================================
\begin{thebibliography}{99}

\bibitem{hollowllm}
C.~Gong, B.~Liu, and M.~Li, ``Hollow-LLM Attack: Computationally Trivial
Weights in Zero-Knowledge Verification of LLM Inference,'' in \emph{IEEE S\&P},
2026, pp.~4339--4355, DOI 10.1109/SP63933.2026.00258.

\bibitem{zkattach}
\need{zk-attach, arXiv:2608.30083 --- 打开原文核实作者后再引}

\bibitem{zkgen}
J.~Lauinger, J.~Ernstberger, and S.~Steinhorst, ``zkGen: Policy-to-Circuit
Transpiler,'' in \emph{IEEE ICBC}, 2024, pp.~620--624,
DOI 10.1109/ICBC59979.2024.10634440.

\bibitem{pco}
\need{PCO, IACR ePrint 2026/994 --- 打开原文核实作者（Dolev/Yung 等）后再引}

\bibitem{zkagent}
L.~Wang, H.~Lou, C.~Li, Y.~Yu, and Y.~Hu, ``zkAgent: Verifiable LLM Agent
Execution via One-Shot Transcript Proofs,'' IACR ePrint 2026/199.

\bibitem{vdr}
\need{Verifiable Document Redacting, Springer, DOI 10.1007/978-3-319-66402-6_20
--- 核实作者与年份。\textbf{绝不能再引 IACR ePrint 2024/813}（那是
Kiraz/Larraia/Vaughan, "How to Redact the Bitcoin Backbone Protocol"）。
见 docs/tdsc-fixlist.md \S A。}

\end{thebibliography}

\end{document}
"   % 两次，第二遍收交叉引用
%  \SKELQUIET 让所有中文标记宏吞掉参数；不带它则需要能排版中文的引擎
%  （xelatex + CJK 字体），本骨架未配置。
% =====================================================================
\documentclass[journal]{IEEEtran}

\usepackage{amsmath,amssymb,amsthm}
\usepackage{booktabs}
\usepackage{array}
\usepackage{graphicx}
\usepackage{xcolor}
\usepackage{hyperref}

% 窄栏里用 p{} 会因两端对齐产生大片字距（编译时刷一堆 Underfull \hbox）。
% 表格正文列一律用 L{}：左对齐、允许自然参差。
\newcolumntype{L}[1]{>{\raggedright\arraybackslash}p{#1}}
\newcommand{\tabfont}{\footnotesize}

\newtheorem{theorem}{Theorem}
\newtheorem{lemma}[theorem]{Lemma}
\newtheorem{proposition}[theorem]{Proposition}
\newtheorem{corollary}[theorem]{Corollary}
\newtheorem{definition}[theorem]{Definition}
\newtheorem{assumption}[theorem]{Assumption}
\newtheorem{remark}[theorem]{Remark}

% --- 工作标记：提交前全部清掉 ---
% 注意：宏名不能叫 \wp —— \wp 已是 LaTeX/amssymb 的 Weierstrass ℘，
% \newcommand{\wp} 会直接报 "Command \wp already defined" 并使后文级联报错。
\ifdefined\SKELQUIET
  \newcommand{\wpmark}[1]{}
  \newcommand{\need}[1]{}
  \newcommand{\fixme}[1]{}
\else
  \newcommand{\wpmark}[1]{\textcolor{red}{\textsuperscript{[WP#1]}}}
  \newcommand{\need}[1]{\textcolor{blue}{[[ NEED: #1 ]]}}
  \newcommand{\fixme}[1]{\textcolor{orange}{[[ FIX: #1 ]]}}
\fi

\begin{document}

% ---------------------------------------------------------------------
% 标题：决策 D5 —— 保留 "Zero-Knowledge" 的前提是 groth16 成为交付档 (WP1)。
%       若 WP1 未打通 groth16，必须删掉标题与摘要里的 zero-knowledge。
% ---------------------------------------------------------------------
\title{Proof-of-Policy: Zero-Knowledge Compliance Certificates\\ for LLM Agent Outputs}
% 备选标题（若 groth16 未打通）："Proof-of-Policy: Succinct Compliance
% Certificates for LLM Agent Outputs"

\author{\need{authors}}
\maketitle

% ---------------------------------------------------------------------
\begin{abstract}
% 决策 E2：内容隐私上界必须出现在摘要，不能埋进 §8.1。
We present \emph{Proof-of-Policy} (PoP), a system that proves an LLM agent's
response $T$ satisfies a publicly declared policy $\pi$, without a trusted
execution environment...
\need{一句完整的贡献陈述（对应 D1 的"分析 x 系统各半"）}

\smallskip
\noindent\textbf{Scope of privacy.} PoP provides \emph{policy}-zero-knowledge:
the certificate reveals neither the policy evaluation trace nor, in
policy-private mode, the policy itself. It does \emph{not} guarantee that $T$
is unrecoverable: as we show in \S\ref{sec:binding}, any publicly recomputable
binding to $T$ admits a re-execution attack whenever the adversary knows the
prompt and the model. The same bit-budget governs the request side: a
certificate is replayable without limit unless the receiver contributes fresh
randomness (Theorem \ref{thm:t3}).
\end{abstract}

\begin{IEEEkeywords}
zero-knowledge proofs, LLM agents, policy compliance, verifiable computation,
accountability
\end{IEEEkeywords}

% =====================================================================
\section{Introduction}
\label{sec:intro}

\need{Contribution list. 必须包含下面这条"不主张"声明：}

\noindent\textbf{What we do \emph{not} claim.}
We do not claim that policy proofs bind nothing: our own system already binds
the response via a verifier-issued nonce (Definition 1). We do not claim to be
the first to observe that a proof can fail to bind the artifact it purports to
describe: the \emph{effort gap} for ZK inference proofs is formalised by
Hollow-LLM \cite{hollowllm}, and the \emph{source-integrity gap} for predicate
proofs is named by zk-attach \cite{zkattach}. Our contribution is to
(i) show that the surviving gap is the \emph{divergent-delivery} one
(Proposition \ref{thm:gap}), (ii) give a quantitative bound on what a public
value reveals, together with the \emph{structural hypothesis} under which that
bound holds---and a demonstration that the hypothesis is necessary, since our
own two public values sit on opposite sides of it (Theorem \ref{thm:t1} and
Remark \ref{rem:betanecessary}), and (iii) give a quantitative price for
closing the gap (Theorem \ref{thm:t3}).

\smallskip
\noindent\textbf{A separation we deliberately do \emph{not} claim.} It is
tempting to promote our designated-verifier interface into a theorem
(``public verification is impossible when $k$ is small''). We give the reason
we do not in Remark \ref{rem:designated}: in the low-$k$ regime hiding fails
for every scheme, designated or not, and the dichotomy has no content of its
own. We state this explicitly because the negated claim is the more attractive
one to write.

\smallskip
\noindent\textbf{One spine, two sides.} Theorems \ref{thm:t1} and
\ref{thm:t3} state the same fact about two different inputs: \emph{a proof
binds only as much as the committed, adversary-non-selectable quantity has bits
for}. The response side is governed by entropy,
$k = -\log \mathsf{Guess}(T \mid Z)$ (Definition \ref{def:k}); the request side
by freshness, $f$. Neither can be bought with more computation.

\fixme{全文禁止出现"比 Proof-of-Guardrail 更强"或任何强弱对比。
Proof-of-Guardrail 与本文正交：它证"计算发生过"，本文证"性质成立"。}

% =====================================================================
\section{Background and Related Work}
\label{sec:related}

\need{五条线，每条都要"它占什么位 / 我们差在哪"两列。见 docs/tdsc-plan.md \S2。}

\begin{table}[t]
\centering
\caption{Positioning against the five prior-art lineages.}
\label{tab:related}
\setlength{\tabcolsep}{3pt}
\begin{tabular}{@{}L{0.26\columnwidth}L{0.42\columnwidth}L{0.26\columnwidth}@{}}
\toprule
Lineage & Representative work & Occupied slot \\
\midrule
Policy $\to$ circuit      & zkGen \cite{zkgen}, Circuitree, Prezta & DSL-to-circuit \\
Agent/LLM compliance      & NiyamAI, Zaguk, PCO \cite{pco}, WITNESS & private-mode kernel \\
Model computation         & DeepProve, zkGPT, zkAgent \cite{zkagent} & inference integrity \\
Attacks / negative results& Hollow-LLM \cite{hollowllm}, zk-attach \cite{zkattach} & naming the gap \\
Non-ZK industry baseline  & @brivora/verify, VAOL & signed audit logs \\
\bottomrule
\end{tabular}
\end{table}

\fixme{引文纪律：除 Hollow-LLM / zkGen / NiyamAI / zkAgent / ePrint 2024/813
五条外，其余引文必须逐条打开原文核实后再进 bibliography。
见 docs/tdsc-fixlist.md \S2 的注记。}

% =====================================================================
\section{System Overview}
\label{sec:system}
% 决策 D1 的"系统半边"：这一节压缩到 ~1.5 页，是 T3(b) 的载体，不是主角。

\need{DSL $\to$ ConstraintSpec $\to$ circuit；双模式；证书 + 锚定；
工具回执链；801 测试 / 15 skip 的工件陈述。}

% =====================================================================
\section{Binding of Policy Proofs}
\label{sec:binding}
% 决策 D1 的"分析半边"：这是论文的核心。

\subsection{The gap}

\begin{definition}[Policy-proof system]
A policy-proof system is a tuple $(\mathsf{Prove}, \mathsf{Ver}, \pi, \mathcal{M})$.
On challenge $n$ from a verifier, $\mathsf{Prove}(\pi, T, n)$ outputs a
certificate $c$ whose public values include
\begin{multline*}
b \;=\; H(\texttt{pop-bind-v1} \,\|\, \mathsf{u32be}(|n|) \,\|\, n \,\|\, T),\\
C = H(T).
\end{multline*}
$\mathsf{Ver}(c, n, T')$ accepts iff the policy relation holds \emph{and}
$b$ recomputed from the received $T'$ matches.
\end{definition}

\fixme{\textbf{H3/H4}：早先写的 ``the relation contains no term referring to the
delivered bytes'' 与仓库事实冲突——公开值里\emph{已有} $b$，验证时由送达的
$T'$ 现场重算（\texttt{circuits/types/src/lib.rs:444}、
\texttt{docs/security-model.md:232}）。命题已按上面的定义改写。见
\texttt{docs/tdsc-fixlist.md} \S H。}

\begin{proposition}[Divergent delivery]
\label{thm:gap}
Definition 1 binds $T'$ \emph{as received by the verifying party}. If the
provider delivers a different string $T'' \ne T'$ to the end user, verification
still succeeds on $T'$. Neither hardware nor a proof can prevent this; only the
end user holding $n$ and recomputing $b$ can.
\end{proposition}
\begin{proof}
$b$ is a function of $(n,T)$ only; the delivery channel to the end user carries
no authenticated term. One line.
\end{proof}

\need{This is the \emph{residual} gap, not ``proofs bind nothing''. Prior work
must be cited here (Hollow-LLM \cite{hollowllm}, zk-attach \cite{zkattach}), and
we state plainly that the naive claim would be false against our own system.}

\subsection{A quantitative binding--hiding tradeoff}
\label{sec:tradeoff}

% ---------------------------------------------------------------------
% T1 : 关键是把"对手已知信息 Z"定对。
%      docs/tdsc-plan.md §1.4 已修正原 §4 的非因果推理。
% ---------------------------------------------------------------------
\begin{definition}[Guessing entropy of a response]
\label{def:k}
For a response distribution $\mathcal{D}$ and adversary side information $Z$,
let
\[
\mathsf{Guess}(T \mid Z) \;\triangleq\; \max_{\mathcal{A}} \Pr[\mathcal{A}(Z) = T],
\]
which equals $\mathbb{E}_z[\max_t \Pr[T = t \mid Z = z]]$, and
\[
k \;\triangleq\; -\log_2 \mathsf{Guess}(T \mid Z),
\]
the maximum ranging over all time-unbounded adversaries.
\end{definition}

\begin{remark}[The definition is operational, and it is the average]
\label{rem:kmeas}
$k$ is \emph{measured}, not assumed: sampling one model on one prompt many
times estimates $\mathsf{Guess}(T \mid Z)$ directly
(\S\ref{sec:eval}), so Theorem \ref{thm:t1} is instantiated with a
number we obtain rather than a number we posit. The definition must average
over $z$ rather than take the pointwise worst case
$k_z = -\log \max_t \Pr[T = t \mid Z = z]$: the proof takes an expectation over
$z$, and $\mathbb{E}_z[2^{-k_z}] = 2^{-k}$ holds only for the average form.
\end{remark}

\begin{assumption}[Adversary side information]
\label{as:Z}
The adversary knows the prompt, the model (possibly its weights), and the
decoding parameters. \emph{This is the assumption that makes the theorem
non-trivial}: under greedy decoding, $k \approx 0$.
\end{assumption}

\begin{theorem}[Binding, and the hypothesis it needs]
\label{thm:t1}
Let $\beta : \mathcal{M} \to \{0,1\}^\lambda$ be public, deterministic and
efficiently computable, and let the certificate publish $\beta(T)$. Let
$\mathcal{A}$ be any adversary, unbounded in time, that is given $Z$ and
evaluates $\beta$ at most $q$ times.
\begin{enumerate}
\item[(H)] \emph{(Randomised public value.)} Suppose
  $\beta(T) = H(\mathsf{dom} \,\|\, \mathsf{u32be}(|n|) \,\|\, n \,\|\, T)$ for a
  random oracle $H$ and a nonce $n$ drawn uniformly and \emph{independently of
  $T$}. Then
  \[
  \Pr\!\left[\mathcal{A}^\beta\!\left(Z, \beta(T)\right) = T\right]
  \;\le\; q \cdot 2^{-k} \;+\; q \cdot 2^{-\lambda}.
  \]
  Consequently hiding with advantage $\varepsilon$ requires
  $k \ge \log q + \log(1/\varepsilon)$.
\item[(L)] \emph{(Lossy public value.)} For an arbitrary deterministic public
  $\beta$ the correct bound is over the \emph{preimage class}:
  \begin{multline*}
  \Pr\!\left[\mathcal{A}^\beta\!\left(Z, \beta(T)\right) = T\right]
  \;\le\; \mathsf{Guess}\!\left(T \mid Z, \beta(T)\right) \\
  \triangleq\; \mathbb{E}_y\!\left[\max_t \Pr[T = t \mid Z, \beta(T) = y]\right].
  \end{multline*}
\end{enumerate}
\end{theorem}

\begin{proof}
\emph{(H).} Fix $z$ and let the adversary's query set be
$Q = \{q_1, \dots, q_q\}$. Decompose by first hit,
\[
\Pr[\exists i : q_i = T] \;=\; \sum_i \Pr\!\left[q_i = T \wedge \neg E_{<i}\right],
\]
where $E_{<i} \triangleq \{\exists j < i : q_j = T\}$.
Conditioned on $\neg E_{<i}$, every oracle answer $a_1,\dots,a_{i-1}$ already
revealed is uniform and independent of $T$ (no queried point equals $T$, and $H$
is a fresh random function), and $\beta(T)$ itself is uniform and independent of
$T$. Hence $q_i$---a deterministic function of $(z, \beta(T), a_{<i})$---is
\emph{independent of $T$}, so
\begin{multline*}
\Pr[q_i = T \mid Z = z]
= \sum_t \Pr[q_i = t]\,\Pr[T = t \mid Z = z] \\
\le \max_t \Pr[T = t \mid Z = z] = 2^{-k_z},
\end{multline*}
giving $\Pr[\exists i : q_i = T \mid z] \le q\,2^{-k_z}$. Taking the expectation
over $z$ and using $\mathbb{E}_z[2^{-k_z}] = \mathsf{Guess}(T \mid Z) = 2^{-k}$
from Definition \ref{def:k} yields $q \cdot 2^{-k}$. An adversary that outputs
an unqueried value wins with probability at most $2^{-k_z} \le q\,2^{-k_z}$,
which the same bound absorbs. The term $q \cdot 2^{-\lambda}$ charges the
probability that a query hits the uniform value $\beta(T)$ by chance.

\emph{(L).} Immediate from Definition \ref{def:k} after conditioning on
$\beta(T)$; no hypothesis on $\beta$ is used.
\end{proof}

\begin{remark}[The hypothesis in (H) is necessary, and our two instances straddle it]
\label{rem:betanecessary}
Without the hypothesis, part (H) is not merely unproven but \emph{false}: take
$\beta = \mathrm{id}$, or $\beta = \mathsf{encode}$ as in
Proposition \ref{prop:semantic}. The adversary outputs $\beta(T)$ and succeeds
with probability $1$, independently of $q$ and $k$. Our own two public values
therefore sit on opposite sides of the hypothesis---$b =
H(\mathsf{dom}\|n\|T)$ satisfies it, $\mathsf{encode}(T)$ does not---and this is
what makes the pair a theorem about \emph{our design} rather than a citation of
a standard bound. An earlier draft of this work stated (H) without its
hypothesis; the counterexample above is why the hypothesis is now explicit.
\end{remark}

\begin{corollary}[Re-execution, not enumeration]
\label{cor:reexec}
Under Assumption \ref{as:Z}, when the decoding is greedy or low-temperature,
$k \approx 0$ and the attack of Theorem \ref{thm:t1} is realised by
\emph{re-running the model}, at the cost of one inference---independent of
$|T|$ and requiring \emph{no} evaluation of $\beta$. Enumeration over the string
space is \emph{not} the relevant attack.
\end{corollary}

\fixme{原 docs/sp1-zk-audit.md §4 写的是"熵(~10^2--10^3 bit)远低于 256 bit，
因此枚举可行"——这是非因果推论（2^100 与 2^256 一样枚举不了；真正的攻击是\textbf{重执行}，
对手的已知信息是 (prompt, 模型, 解码参数)）。
2026-10-12 已全库改正，docs/tdsc-fixlist.md §G 的 G1--G6 逐条结案。}

\begin{remark}[A design observation, not a separation]
\label{rem:designated}
Our verification interface is \emph{designated}: binding is checked against a
value $b = H(\mathsf{dom}\|n\|T)$ in which a verifier-issued nonce $n$ appears.
It is tempting to promote this into a separation (``public verification is
impossible when $k$ is small''). \textbf{We claim no such theorem, and
Corollary \ref{cor:reexec} is why.} When $k \approx 0$, hiding fails for
\emph{every} scheme, designated or not, because the adversary re-runs the
model; when $k$ is large, a zero-knowledge proof of ``$\exists\, T : b =
H(\mathsf{dom}\|n\|T) \wedge J(\pi, T) = 1$'' is simultaneously publicly
verifiable and hiding. The dichotomy is thus governed by $k$ alone and carries
no independent content. What does survive---and what
\S\ref{sec:wp4-tension}
develops---is that \emph{third-party auditing of delivery} forces disclosure of
$T$, which is an informational remark rather than a theorem. We record the
demotion here because an earlier draft asserted the separation as a theorem.
\end{remark}

\need{\textbf{This is retrospective justification, not a novel construction.}
The system publishes $b = H(\mathsf{dom}\|n\|T)$ with a verifier-issued nonce
$n$, and Remark \ref{rem:designated} explains the shape of that choice. The
security model does not pin down whether $n$ is public; making it public
forfeits hiding, and making it verifier-held is what makes verification
designated.}

\need{\textbf{实验内容}：实测 $k$。对一组真实 prompt、同一模型多次采样，
估 $\mathsf{Guess}(T\mid Z)$（Definition \ref{def:k}），代入
Theorem \ref{thm:t1}(H) 得阈值。
据先例扫查，\emph{没有人测过这个量}——它本身是一块贡献。}

\begin{proposition}[The lossy instance of Theorem \ref{thm:t1}(L)]
\label{prop:semantic}
For the semantic rule class, the public values contain $\beta(T) =
\mathsf{encode}(T)$, a content-preserving feature vector of the response. This
$\beta$ is deterministic and injective on the features but not on strings, so
it violates the hypothesis of Theorem \ref{thm:t1}(H) and instantiates part
(L): the binding force is exactly
$\mathsf{Guess}(T \mid Z, \mathsf{encode}(T))$, i.e.\ the guessing probability
\emph{within the preimage class} of the published vector. Here $k$ need not be
argued from perplexity at all: the adversary's search space is that preimage,
which is dramatically smaller than the string space. The relevant asymmetry is
not ``low entropy'' but a \emph{lossy public encoding}.
\end{proposition}
\begin{proof}
Immediate. The measurement is direct: publish $\mathsf{encode}(T)$ for real
responses and measure how much of $T$ is recoverable.
\end{proof}

\need{This proposition is the paper's cheapest experiment and its strongest
anchor for Theorem \ref{thm:t1}---it is measurable in a day, and it makes the
bound concrete rather than cited. Source: \texttt{docs/design-semantic-rules.md}
lines 195--201, where the repository itself already names it as a second
instance.}

\subsection{Quantitative request binding}
\label{sec:t3}
% 决策 D2（第二版，2026-10-12）：原「闭合三分法」定理降级为 Remark（分类），
% 取代它的是下面的定量界 T3′。理由见 docs/tdsc-plan.md §1.6(b)：
% 三分法要么假（理想功能模型下能构造出第四条路），要么琐碎（把三条路写成定义）。

\begin{theorem}[Selective-presentation advantage]
\label{thm:t3}
Fix a policy $\pi$ and a request distribution. A selective-presentation
adversary $\mathcal{B}$ may, \emph{before} seeing the receiver's challenge,
compute a set $S$ of accepted transcripts. The receiver then contributes $f$
fresh, uniformly random bits $r$ that enter the proven statement. Then
\[
\Pr\!\left[\mathcal{B}\ \text{is accepted on } r\right] \;\le\; |S| \cdot 2^{-f}.
\]
If instead the receiver contributes no fresh randomness ($f = 0$; e.g.\ an
empty nonce), the advantage is $1$, and \emph{does not decay in any parameter
of the system}.
\end{theorem}
\begin{proof}
\need{主论证：$r$ 进入 $b = H(\text{dom} \| r \| T)$，故 $S$ 里任何一个
预先算好的证书，只有在 $r$ 恰好落进它已覆盖的挑战集合时才被接受；
$f$ 个新鲜均匀比特给出 $2^{-f}$ 的对撞上界，再对 $S$ 取并界。
$f = 0$ 时 $S$ 中每个元素都是永久有效的重放。}

\need{\textbf{代码依据}：\texttt{circuits/types/src/lib.rs}:950--995 已写明
“nonce 为空是合法的……不提供任何重放防护”。定理的 $f = 0$ 分支不是假想，
是本系统当前 API 的\emph{合法输入}。}
\end{proof}

\need{这条定理与 Theorem \ref{thm:t1} \textbf{同构}，而正是这一点使它值得写：
两条说的是同一件事——\emph{一个证明只能绑住提交进去的、对手无法预先选择的
那部分比特所能支撑的量}。Theorem \ref{thm:t1} 管响应侧
（$k = -\log_2 \mathsf{Guess}(T \mid Z)$，见 Definition \ref{def:k}），Theorem \ref{thm:t3} 管请求侧（新鲜度 $f$）。
这条统一主线写在 \S\ref{sec:intro} 的贡献列表里，是全文的理论骨架。}

\begin{remark}[A taxonomy, not a trichotomy]
\label{rem:taxonomy}
Every policy-proof deployment we surveyed closes the divergent-delivery gap of
Proposition \ref{thm:gap} in one of three ways:
(a) root trust in hardware (attested enclaves);
(b) have the \emph{end user} contribute a secret to the proven statement (the
user issues the nonce); or
(c) gate delivery on successful verification (\emph{prove-before-deliver}).
\textbf{We offer this as a classification, not as an exhaustiveness theorem}:
we neither prove that a fourth family cannot exist, nor that the three are
disjoint. zk-attach \cite{zkattach} walks (a); NiyamAI and ZKAP walk (c); PoP
walks (b).
\end{remark}

\need{\textbf{Posture, stated plainly.} Our system already carries
\texttt{response\_binding}, but $n$ is issued by the \emph{verifier}. Moving it
to the \emph{end user}---making $f > 0$ the receiver's contribution---closes
the divergent-delivery branch of Proposition \ref{thm:gap}. This is a small
protocol amendment, not a new system, and the paper must say exactly that. What
Theorem \ref{thm:t3} buys is the \emph{price}: with $f = 0$ there is no closure
at any cost.}

% =====================================================================
\section{Policy-Private Compliance}
\label{sec:private-policy}
% 决策 D3：这是"透明模式 ZK 冗余"这个批评的唯一硬解。
% 2026-10-12 读代码后的定性：这不是"造一个新证明系统"，是一次\emph{披露策略}
% 的重构——策略文本在架构上已经是证人。见 docs/tdsc-plan.md WP4.1。

\need{动机：对确定性规则 + 公开响应，签名哈希链就够了。唯一无法用签名
替代的场景是策略本身要保密。}

\subsection{The observation that reframes this work package}
\label{sec:wp4-finding}

\texttt{spec\_canonical} is already a \emph{guest input}---a witness---not a
public value: only
$\mathsf{policy\_hash} = \mathsf{SHA256}(\texttt{spec\_canonical})$
enters the public values (\texttt{circuits/types/src/lib.rs}:1--16, 415--416,
918, 1003--1004). So ``do not reveal the policy text'' is
\emph{architecturally already true}; what leaks is everything \emph{derived}
from the policy. WP4 is accordingly a \emph{disclosure-policy refactor}, not a
new proof system, and must be written that way---as a systematic account of
what a certificate leaks about the policy it certifies.

\begin{table}[t]
\centering
\caption{What a policy-private certificate leaks today, and what closing each
leak costs.}
\label{tab:wp4}
\begin{tabular}{@{}L{0.04\linewidth}L{0.46\linewidth}L{0.36\linewidth}@{}}
\toprule
\# & Leak & Difficulty / dependency \\
\midrule
1 & The policy text itself & \textbf{Already closed} (it is a witness). \\
2 & Witness hiding in the proof system & Medium; \emph{inherited from WP1}:
    only \texttt{groth16}/\texttt{plonk} is witness-hiding, per
    \texttt{docs/\allowbreak sp1-zk-audit.md}. \\
3 & \texttt{PrivateViolation}\allowbreak\{\texttt{rule},\allowbreak\texttt{kind}\} names the rule that
    fired & Medium; needs a coarsened violation class plus a commitment to it. \\
4 & Evidence opening reveals the rule
    (\texttt{policydsl/\allowbreak privacy/\allowbreak commit.py}:319) & Medium; disclose the
    \emph{existence} of a violation, not its identity. \\
5 & The \texttt{delegated} flag exposes policy structure (which rules were
    routed to the ezkl path) & Medium; easiest to overlook, cheapest to fix. \\
6 & Semantic rules publish $\beta(T)=\mathsf{encode}(T)$
    (Proposition \ref{prop:semantic}) & Medium--high; either bound the
    preimage, or disable this rule class in policy-private mode. \\
\bottomrule
\end{tabular}
\end{table}

\need{Honest structure of the table, and it should be stated as such:
sub-problem 1 is a one-paragraph re-reading of the existing code, \emph{not} a
contribution; 2 is inherited wholesale from WP1; 3--5 are engineering with one
real formal core (disclosure granularity); 6 is where the work is---and it
collides with Theorem \ref{thm:t1}, because one cannot hide the policy while
publishing a lossy encoding of the response that the semantic rules require.}

\subsection{The tension worth writing up: policy privacy vs.\ accountability}
\label{sec:wp4-tension}

Lead the section with the tension, not with the cryptography: it determines
whether the mode means anything at all.

\paragraph{The prerequisite, stated as a prerequisite.} Let
$H = \mathsf{SHA256}(\pi)$. On its own $H$ is \emph{inert}. A prover can always
emit a self-consistent certificate for any policy $\pi'$ it prefers, simply by
placing $H' = \mathsf{SHA256}(\pi')$ in that certificate. \emph{Secrecy of
$\pi$ is what creates the gap}, and no cryptographic mechanism can close it:
closing it would require showing $\pi$ to the verifier, which is precisely what
the mode exists to avoid. What closes it is an \emph{issuing act}, so we state
it as an explicit prerequisite of the mode:

\begin{description}
  \item[(P-PP)] An \emph{issuing authority} signs $\pi$ and registers $H$ in a
  record readable by third parties. A policy is \emph{accountable} iff its
  \texttt{policy\_hash} can be traced to such a record.
\end{description}

Without (P-PP) the mode is not merely weaker---the verifier's conclusion is
empty. This is the same assumption as ``honest policy author'' in
\texttt{docs/security-model.md}~\S0.2; what is new here is scoping it to the
verification interface, which had not been done, and which a referee will
otherwise use to dismiss the whole mode in one sentence.

\paragraph{Three parties.} ``Three'' means \emph{authority / operator / third
party}; an auditor is a named instance of the third party, not a fourth trust
anchor.

\begin{table}[t]
\centering
\caption{The three-party model of policy-private mode.}
\label{tab:wp4-parties}
\setlength{\tabcolsep}{3pt}
\begin{tabular}{@{}L{0.15\columnwidth}L{0.13\columnwidth}L{0.64\columnwidth}@{}}
\toprule
Party & Honest? & Holds; wants \\
\midrule
Authority $R$ & \textbf{Yes}, by (P-PP) & Holds $\pi$ and the issuing record;
wants $\pi$ actually enforced, and wants to be able to check that it was. \\
Operator $O$ & \textbf{No} (the adversary) & Holds $\pi$, $T$, the proving key;
wants $\pi$ hidden from third parties while still being able to certify. \\
Third party $V$ & Yes & Holds only public values---\texttt{policy\_\allowbreak hash},
\texttt{circuit\_\allowbreak hash}, \texttt{response\_\allowbreak commitment}, \texttt{passed}, the
proof; wants to confirm that \emph{some issued} policy passed, and must
\emph{not} learn $\pi$. \\
\emph{(Auditor $V^{*}$)} & Yes & As $V$, plus $\pi$ once $O$ discloses it;
wants ex-post accountability. \\
\bottomrule
\end{tabular}
\end{table}

Two boundaries must be stated with the model. First, secrecy of $\pi$ comes
from a \emph{data-flow fact}---$\pi$ never enters the public values---not from
zero-knowledge; any rule that writes a fragment of $\pi$ into the public values
(the leaks in Table~\ref{tab:wp4}, items~3 and~5) breaks it on that fragment.
Second, accountability for $V^{*}$ is \emph{voluntary}: $V^{*}$ can check only
when $O$ chooses to disclose $\pi$. The mode offers no forced disclosure; that
would need the release-gating family discussed in
\texttt{docs/tdsc-plan.md}~\S1.6(b), which we do not claim.

\paragraph{Verifier capability table.} This table is the contribution; the
cryptography is downstream of it, and the table survives even if the
cryptography only partially lands. \emph{It must appear in the paper as
written}, because it is what stops \texttt{policy\_hash} from being read as
``compliance was proved''.

\begin{table}[t]
\centering
\caption{What each verifier may and may not conclude in policy-private mode.}
\label{tab:wp4-capability}
\footnotesize
\setlength{\tabcolsep}{3pt}
\begin{tabular}{@{}L{0.14\columnwidth}L{0.31\columnwidth}L{0.30\columnwidth}L{0.17\columnwidth}@{}}
\toprule
Verifier & \textbf{May} conclude & \textbf{May not} conclude & Rests on \\
\midrule
Third party $V$ & (i) the commitment binds $T$
(Definition~\ref{def:k}); (ii) \emph{some} policy with hash $H$ passed;
(iii) $H$ appears in the issuing record & (iv) the \emph{content} of $\pi$;
(v) how many rules $\pi$ has, or of which classes; (vi) ``the operator
enforces the regulation \emph{I} care about''---only ``some issued
regulation'' & (P-PP); A2; the $k$ of T1 \\
End user (also receiver) & As $V$, and additionally: the proved $T$ is the
string it received, since its nonce is bound (T3$'$, requires $f>0$) & Only
the string it received; \emph{not} what was delivered to another downstream
agent & $f>0$ \\
Authority $R$ & Everything; it holds $\pi$ and can recompute $H$ & --- & --- \\
Auditor $V^{*}$ (after disclosure) & (vii) $\mathsf{SHA256}(\pi) =
\texttt{policy\_\allowbreak hash}$, i.e.\ what ran was the issued $\pi$;
(viii) the disclosed violation record is consistent & (ix) the undisclosed
part of $T$ & voluntary disclosure \\
On-chain verifier & (i)--(iii) only; it hashes and verifies signatures, and
does not interpret policy semantics & (iv)--(ix) & (P-PP) \\
\bottomrule
\end{tabular}
\end{table}

The $V$ row is the whole reason this subsection exists: $V$ does not obtain
``compliance'', it obtains ``some issued policy was satisfied''. Conflating the
two is the sentence a referee would use to kill the mode.

\need{Remaining work on this subsection: (a) confirm the (P-PP) framing with
the collaborator, and choose its instantiation (an on-chain registry of issued
policy hashes vs.\ a signed manifest); (b) add the prior-art citations below;
(c) decide whether the capability table merges with Table~\ref{tab:wp4} in the
submission, or stays separate (currently separate, on the grounds that one
table is about leakage and the other about verifier conclusions).}

\subsection{Expected results, stated in advance}
\label{sec:wp4-expect}

\need{Calibrated expectations, so that a partial outcome is a publishable
result rather than a failure:}
\begin{itemize}
  \item Sub-problems 1, 2, 6: high confidence ($\gtrsim 80\%$). 1 is already
        true of the code; 2 follows from WP1; 6 is a design restriction.
  \item Sub-problems 3, 4, 5: likely ($\sim 70\%$).
  \item ``\emph{Zero}-knowledge violation evidence'' in the strong sense:
        only $50$--$60\%$. The boolean ``some violation occurred'' is easy; the
        hard part is \emph{granularity}---with $3$--$8$ rules, a coarsened
        violation class can itself identify the rule. \textbf{The likely best
        result is a granularity/actionability/leakage tradeoff curve, not full
        hiding}---and that curve is a contribution in its own right.
  \item \S\ref{sec:wp4-tension}, the tension and the capability table:
        $\sim 90\%$. This is precisely why the section leads with it.
\end{itemize}

\need{Prior-art discipline, and this is the highest-risk place in the paper ---
the collaborator's own field. \emph{Policy hiding is not a new primitive}:
hidden credentials, secret handshake, and XACML privacy in access control, and
attribute-hiding / fuzzy predicate encryption, are all the same family. Cite
that line explicitly and do not claim primacy. What is new is narrow: (a) the
\emph{application} to LLM output compliance together with selective disclosure
of violation type and location, and (b) the accountability tension of
\S\ref{sec:wp4-tension}.}

% =====================================================================
\section{Composition: The Constructive Half of the End-User-Secret Family}
\label{sec:composition}
% 决策 D4：双载体（ezkl + 真 ONNX 主臂；DeepProve 对照臂）。

\fixme{关键更正：组合本身不闭合缺口。只有当模型证明绑定到\emph{用户提交的
输入} $x$ 时，$M(x) = T \wedge \pi(T)$ 才真的闭合"送错字节"——而那恰好是
Theorem \ref{thm:t3} 要求的 $f > 0$：用户提交的 $x$ 就是那 $f$ 个新鲜比特。}

\need{组合契约的定义（模型证明须暴露什么，策略证明才能绑上）；
真模型规模（$\ge 10^6$ 参数）；成本数据。}

% =====================================================================
\section{Availability: The Unanchored-Pattern Cost Attack}
\label{sec:dos}
% WP6。

\begin{theorem}[Quadratic cost inflation]
\need{对含未锚定模式的正则，控制响应的对手可强制 $\Omega(n^2)$ 证明工作量。}
\end{theorem}

\need{缓解：Pike VM（已有）+ 响应长度承诺 + 分块匹配。
zk-regex 有 DFA 限制但没有 zkVM 场景下的 DoS 分析——这是新的。}

% =====================================================================
\section{Coverage: Can a Provider Simply Withhold Certificates?}
\label{sec:coverage}
% WP5。诚实预期：无可信硬件时可能无解 —— 那就写成 open problem。

% =====================================================================
\section{Evaluation}
\label{sec:eval}

\need{矩阵：规则 \{1,2,3,4,6,8\} $\times$ 长度 \{200, 2k, 10k, 50k\}
$\times$ 模式 \{transparent, private, private-policy\}
$\times$ 证明档 \{core, groth16\}。
骨架见 bench/bench_scale.py。}

\subsection{Why zero-knowledge? A decision map}
% WP9：非 ZK 基线 + 决策图。

\need{签名 Merkle 日志对照臂（brivora/VAOL 风格）。
结论：ZK 只在 \{私有响应, 私有策略\} 两个角上有正收益。
把"透明模式 ZK 冗余"从批评变成有刻度的设计结论。}

% =====================================================================
\section{Limitations}
\label{sec:limitations}

\need{必须写进去的：下列五条。每条都要在正文里展开，不是罗列。}
\begin{itemize}
  \item Content privacy here is \emph{commitment} privacy, not
        irrecoverability of content (Corollary \ref{cor:reexec}).
  \item The zero-knowledge property of the default proof mode:
        \fixme{SP1's published security model states explicitly that
        individual STARK proofs do not currently satisfy the
        zero-knowledge property. If groth16 does not become the delivery
        mode, the word ``zero-knowledge'' must leave the title.}
  \item Binding the composition's input requires the user to be online (or to
        hold a key). This is the price of Remark \ref{rem:taxonomy}(b), and
        the concrete form of $f > 0$ in Theorem \ref{thm:t3}.
  \item The coverage problem (\S\ref{sec:coverage}) may be unsolvable without
        trusted hardware.
  \item Three-party collusion is undetectable at the signature layer
        (\texttt{docs/security-model.md}).
\end{itemize}

% =====================================================================
\section{Conclusion}

% =====================================================================
\appendix
\section{Proofs}
\need{\textbf{状态（2026-10-12 自检后）}：
Theorem \ref{thm:t1}(H) 与 Theorem \ref{thm:t3} 的证明\textbf{已写在正文}，
此处不再重复。三分法已降级为分类 Remark \ref{rem:taxonomy}，无需证明。
Theorem \ref{thm:t1}(L) 无闭式，只给测量口径（\S\ref{sec:eval}）。
原 T2「分离定理」已\textbf{停用}（见 Remark \ref{rem:designated}）——
本节因此只剩一条待补事项，见下。}

\need{\textbf{待补}：把 Theorem \ref{thm:t1}(H) 的随机预言机模型换成
标准模型（或明确声明本文只在 ROM 下主张）。当前证明依赖 $H$ 是新鲜随机函数
这一条；实践中的 SHA-256 只给计算意义下的近似，故正文措辞必须写成
\emph{heuristic}，不得写成信息论界。}

% =====================================================================
\begin{thebibliography}{99}

\bibitem{hollowllm}
C.~Gong, B.~Liu, and M.~Li, ``Hollow-LLM Attack: Computationally Trivial
Weights in Zero-Knowledge Verification of LLM Inference,'' in \emph{IEEE S\&P},
2026, pp.~4339--4355, DOI 10.1109/SP63933.2026.00258.

\bibitem{zkattach}
\need{zk-attach, arXiv:2608.30083 --- 打开原文核实作者后再引}

\bibitem{zkgen}
J.~Lauinger, J.~Ernstberger, and S.~Steinhorst, ``zkGen: Policy-to-Circuit
Transpiler,'' in \emph{IEEE ICBC}, 2024, pp.~620--624,
DOI 10.1109/ICBC59979.2024.10634440.

\bibitem{pco}
\need{PCO, IACR ePrint 2026/994 --- 打开原文核实作者（Dolev/Yung 等）后再引}

\bibitem{zkagent}
L.~Wang, H.~Lou, C.~Li, Y.~Yu, and Y.~Hu, ``zkAgent: Verifiable LLM Agent
Execution via One-Shot Transcript Proofs,'' IACR ePrint 2026/199.

\bibitem{vdr}
\need{Verifiable Document Redacting, Springer, DOI 10.1007/978-3-319-66402-6_20
--- 核实作者与年份。\textbf{绝不能再引 IACR ePrint 2024/813}（那是
Kiraz/Larraia/Vaughan, "How to Redact the Bitcoin Backbone Protocol"）。
见 docs/tdsc-fixlist.md \S A。}

\end{thebibliography}

\end{document}
